\documentclass[11pt,a4paper]{article}

\newcommand{\CourseCode}{CSAI2017P: Applied Machine Learning}
\newcommand{\AssignmentName}{A1 -- Theory Assignment 1}

% Shared settings for Applied Machine Learning lab assignments and marking schemes.
% Layered on the house notes preamble; keep free of \documentclass.
%
% Callers must define \CourseCode, \AssignmentName before \begin{document}.

% Depth-agnostic locate of the house notes preamble: briefs compile from
% public/lab-assignments/LAxx/latex/ and marking schemes from
% private/solutions/lab-assignments/LAxx/latex/, which sit at different depths.
\IfFileExists{../../../../public/_shared/notes-common.tex}{%
  % Shared notes settings for Applied Machine Learning lectures (UPES).
% Keep this file free of \documentclass to allow reuse via \input.

\usepackage[utf8]{inputenc}
\usepackage[T1]{fontenc}
\usepackage{lmodern}          % scalable T1 fonts; without this pdflatex falls back
\input{glyphtounicode}        % to bitmapped EC Type 3 and the PDFs are neither
\pdfgentounicode=1            % searchable nor screen-reader friendly
\usepackage[a4paper,margin=1in]{geometry}
\usepackage{hyperref}
\usepackage{xurl}
\usepackage{booktabs}
\usepackage{amsmath}
\usepackage{amssymb}
\usepackage{listings}
\usepackage{xcolor}
\usepackage{enumitem}
\usepackage{parskip}

% Code listings (Python-oriented for this subject).
\lstset{
  language=Python,
  basicstyle=\ttfamily\small,
  keywordstyle=\color{blue},
  commentstyle=\color{gray},
  stringstyle=\color{teal},
  showstringspaces=false,
  columns=fullflexible,
  frame=single,
  framerule=0.2pt,
  breaklines=true
}

\setlist[itemize]{topsep=0.3em,itemsep=0.2em}
\setlist[enumerate]{topsep=0.3em,itemsep=0.2em}

% Simple per-section source line.
\newcommand{\notesources}[1]{\par\noindent{\small\color{gray}\textit{Sources:} \sloppy #1\par}}

% Unit-wise source strings for this subject (used by slides + notes).
% Path is relative to the lecture `latex/` folder (the compilation CWD).
\IfFileExists{../../../_shared/aml-sources.tex}{%
  % Source strings for Applied Machine Learning (CSAI2017P) — used by slides + notes.
% Prescribed texts and reference books are taken from the course syllabus.
% Support texts are the author-hosted open-access editions:
%   statlearning.com            (ISL with Python)
%   probml.github.io/pml-book   (Murphy, Probabilistic Machine Learning)
%   d2l.ai                      (Dive into Deep Learning)
%   mml-book.github.io          (Mathematics for Machine Learning)

\newcommand{\AMLRepoURL}{https://github.com/tali7c/Applied-Machine-Learning}

% --- prescribed -------------------------------------------------------------
\newcommand{\AMLPrescribed}{%
M\"uller and Guido, \textit{Introduction to Machine Learning with Python}, O'Reilly, 2016.;%
Bishop, \textit{Pattern Recognition and Machine Learning}, Springer, 2006.;%
Murphy, \textit{Machine Learning: A Probabilistic Perspective}, MIT Press, 2012.;%
Han, Kamber and Pei, \textit{Data Mining: Concepts and Techniques} (3e), Morgan Kaufmann, 2012.%
}

% --- unit-wise ---------------------------------------------------------------
\newcommand{\AMLUnitISources}{%
M\"uller and Guido, \textit{Introduction to Machine Learning with Python}, Ch. 1.;%
Bishop, \textit{Pattern Recognition and Machine Learning}, Ch. 1.;%
Murphy, \textit{Probabilistic Machine Learning: An Introduction}, MIT Press, 2022, Ch. 1.;%
Mitchell, \textit{Machine Learning}, McGraw Hill, 1997, Ch. 1.;%
Scikit-learn documentation, accessed 2026-08-04.%
}

\newcommand{\AMLUnitIISources}{%
Bishop, \textit{Pattern Recognition and Machine Learning}, \S1.5, \S4.3, \S7.1.;%
Zhang, Lipton, Li and Smola, \textit{Dive into Deep Learning}, \S3.4.;%
Shalev-Shwartz and Ben-David, \textit{Understanding Machine Learning}, Ch. 15.;%
Murphy, \textit{Probabilistic Machine Learning: An Introduction}, Ch. 5.%
}

\newcommand{\AMLUnitIIISources}{%
Zhang, Lipton, Li and Smola, \textit{Dive into Deep Learning}, Ch. 12 (Optimization Algorithms).;%
Deisenroth, Faisal and Ong, \textit{Mathematics for Machine Learning}, Ch. 7.;%
Kingma and Ba, ``Adam: A Method for Stochastic Optimization,'' ICLR 2015.;%
Ruder, ``An overview of gradient descent optimization algorithms,'' arXiv:1609.04747, 2016.%
}

\newcommand{\AMLUnitIVSources}{%
Han, Kamber and Pei, \textit{Data Mining: Concepts and Techniques} (3e), Ch. 3.;%
M\"uller and Guido, \textit{Introduction to Machine Learning with Python}, Ch. 3--4.;%
James, Witten, Hastie, Tibshirani and Taylor, \textit{An Introduction to Statistical Learning with Python}, Ch. 5--6, 12.;%
Scikit-learn documentation (preprocessing, model selection), accessed 2026-08-04.%
}

\newcommand{\AMLUnitVSources}{%
James et al., \textit{An Introduction to Statistical Learning with Python}, Ch. 3--6.;%
Bishop, \textit{Pattern Recognition and Machine Learning}, Ch. 3.;%
Deisenroth, Faisal and Ong, \textit{Mathematics for Machine Learning}, Ch. 9.;%
M\"uller and Guido, \textit{Introduction to Machine Learning with Python}, \S2.3.%
}

\newcommand{\AMLUnitVISources}{%
Han, Kamber and Pei, \textit{Data Mining: Concepts and Techniques} (3e), Ch. 8--9.;%
James et al., \textit{An Introduction to Statistical Learning with Python}, Ch. 8--9.;%
Bishop, \textit{Pattern Recognition and Machine Learning}, Ch. 4--5, 7, 14.;%
Zhang et al., \textit{Dive into Deep Learning}, Ch. 5.%
}

\newcommand{\AMLUnitVIISources}{%
Han, Kamber and Pei, \textit{Data Mining: Concepts and Techniques} (3e), Ch. 10.;%
James et al., \textit{An Introduction to Statistical Learning with Python}, Ch. 12.;%
M\"uller and Guido, \textit{Introduction to Machine Learning with Python}, \S3.5.;%
Ester, Kriegel, Sander and Xu, ``A density-based algorithm for discovering clusters (DBSCAN),'' KDD 1996.%
}

\newcommand{\AMLLabSources}{%
M\"uller and Guido, \textit{Introduction to Machine Learning with Python}, companion notebooks.;%
McKinney, \textit{Python for Data Analysis} (3e), O'Reilly, 2022.;%
Scikit-learn user guide, accessed 2026-08-04.;%
Matplotlib and Seaborn documentation, accessed 2026-08-04.%
}
%
}{%
  % Faculty copies live under `private/` so their compile CWD is different.
  \IfFileExists{../../../../public/_shared/aml-sources.tex}{%
    % Source strings for Applied Machine Learning (CSAI2017P) — used by slides + notes.
% Prescribed texts and reference books are taken from the course syllabus.
% Support texts are the author-hosted open-access editions:
%   statlearning.com            (ISL with Python)
%   probml.github.io/pml-book   (Murphy, Probabilistic Machine Learning)
%   d2l.ai                      (Dive into Deep Learning)
%   mml-book.github.io          (Mathematics for Machine Learning)

\newcommand{\AMLRepoURL}{https://github.com/tali7c/Applied-Machine-Learning}

% --- prescribed -------------------------------------------------------------
\newcommand{\AMLPrescribed}{%
M\"uller and Guido, \textit{Introduction to Machine Learning with Python}, O'Reilly, 2016.;%
Bishop, \textit{Pattern Recognition and Machine Learning}, Springer, 2006.;%
Murphy, \textit{Machine Learning: A Probabilistic Perspective}, MIT Press, 2012.;%
Han, Kamber and Pei, \textit{Data Mining: Concepts and Techniques} (3e), Morgan Kaufmann, 2012.%
}

% --- unit-wise ---------------------------------------------------------------
\newcommand{\AMLUnitISources}{%
M\"uller and Guido, \textit{Introduction to Machine Learning with Python}, Ch. 1.;%
Bishop, \textit{Pattern Recognition and Machine Learning}, Ch. 1.;%
Murphy, \textit{Probabilistic Machine Learning: An Introduction}, MIT Press, 2022, Ch. 1.;%
Mitchell, \textit{Machine Learning}, McGraw Hill, 1997, Ch. 1.;%
Scikit-learn documentation, accessed 2026-08-04.%
}

\newcommand{\AMLUnitIISources}{%
Bishop, \textit{Pattern Recognition and Machine Learning}, \S1.5, \S4.3, \S7.1.;%
Zhang, Lipton, Li and Smola, \textit{Dive into Deep Learning}, \S3.4.;%
Shalev-Shwartz and Ben-David, \textit{Understanding Machine Learning}, Ch. 15.;%
Murphy, \textit{Probabilistic Machine Learning: An Introduction}, Ch. 5.%
}

\newcommand{\AMLUnitIIISources}{%
Zhang, Lipton, Li and Smola, \textit{Dive into Deep Learning}, Ch. 12 (Optimization Algorithms).;%
Deisenroth, Faisal and Ong, \textit{Mathematics for Machine Learning}, Ch. 7.;%
Kingma and Ba, ``Adam: A Method for Stochastic Optimization,'' ICLR 2015.;%
Ruder, ``An overview of gradient descent optimization algorithms,'' arXiv:1609.04747, 2016.%
}

\newcommand{\AMLUnitIVSources}{%
Han, Kamber and Pei, \textit{Data Mining: Concepts and Techniques} (3e), Ch. 3.;%
M\"uller and Guido, \textit{Introduction to Machine Learning with Python}, Ch. 3--4.;%
James, Witten, Hastie, Tibshirani and Taylor, \textit{An Introduction to Statistical Learning with Python}, Ch. 5--6, 12.;%
Scikit-learn documentation (preprocessing, model selection), accessed 2026-08-04.%
}

\newcommand{\AMLUnitVSources}{%
James et al., \textit{An Introduction to Statistical Learning with Python}, Ch. 3--6.;%
Bishop, \textit{Pattern Recognition and Machine Learning}, Ch. 3.;%
Deisenroth, Faisal and Ong, \textit{Mathematics for Machine Learning}, Ch. 9.;%
M\"uller and Guido, \textit{Introduction to Machine Learning with Python}, \S2.3.%
}

\newcommand{\AMLUnitVISources}{%
Han, Kamber and Pei, \textit{Data Mining: Concepts and Techniques} (3e), Ch. 8--9.;%
James et al., \textit{An Introduction to Statistical Learning with Python}, Ch. 8--9.;%
Bishop, \textit{Pattern Recognition and Machine Learning}, Ch. 4--5, 7, 14.;%
Zhang et al., \textit{Dive into Deep Learning}, Ch. 5.%
}

\newcommand{\AMLUnitVIISources}{%
Han, Kamber and Pei, \textit{Data Mining: Concepts and Techniques} (3e), Ch. 10.;%
James et al., \textit{An Introduction to Statistical Learning with Python}, Ch. 12.;%
M\"uller and Guido, \textit{Introduction to Machine Learning with Python}, \S3.5.;%
Ester, Kriegel, Sander and Xu, ``A density-based algorithm for discovering clusters (DBSCAN),'' KDD 1996.%
}

\newcommand{\AMLLabSources}{%
M\"uller and Guido, \textit{Introduction to Machine Learning with Python}, companion notebooks.;%
McKinney, \textit{Python for Data Analysis} (3e), O'Reilly, 2022.;%
Scikit-learn user guide, accessed 2026-08-04.;%
Matplotlib and Seaborn documentation, accessed 2026-08-04.%
}
%
  }{%
    % Question papers live at `private/<family>/<instrument>/`.
    \IfFileExists{../../../public/_shared/aml-sources.tex}{%
      % Source strings for Applied Machine Learning (CSAI2017P) — used by slides + notes.
% Prescribed texts and reference books are taken from the course syllabus.
% Support texts are the author-hosted open-access editions:
%   statlearning.com            (ISL with Python)
%   probml.github.io/pml-book   (Murphy, Probabilistic Machine Learning)
%   d2l.ai                      (Dive into Deep Learning)
%   mml-book.github.io          (Mathematics for Machine Learning)

\newcommand{\AMLRepoURL}{https://github.com/tali7c/Applied-Machine-Learning}

% --- prescribed -------------------------------------------------------------
\newcommand{\AMLPrescribed}{%
M\"uller and Guido, \textit{Introduction to Machine Learning with Python}, O'Reilly, 2016.;%
Bishop, \textit{Pattern Recognition and Machine Learning}, Springer, 2006.;%
Murphy, \textit{Machine Learning: A Probabilistic Perspective}, MIT Press, 2012.;%
Han, Kamber and Pei, \textit{Data Mining: Concepts and Techniques} (3e), Morgan Kaufmann, 2012.%
}

% --- unit-wise ---------------------------------------------------------------
\newcommand{\AMLUnitISources}{%
M\"uller and Guido, \textit{Introduction to Machine Learning with Python}, Ch. 1.;%
Bishop, \textit{Pattern Recognition and Machine Learning}, Ch. 1.;%
Murphy, \textit{Probabilistic Machine Learning: An Introduction}, MIT Press, 2022, Ch. 1.;%
Mitchell, \textit{Machine Learning}, McGraw Hill, 1997, Ch. 1.;%
Scikit-learn documentation, accessed 2026-08-04.%
}

\newcommand{\AMLUnitIISources}{%
Bishop, \textit{Pattern Recognition and Machine Learning}, \S1.5, \S4.3, \S7.1.;%
Zhang, Lipton, Li and Smola, \textit{Dive into Deep Learning}, \S3.4.;%
Shalev-Shwartz and Ben-David, \textit{Understanding Machine Learning}, Ch. 15.;%
Murphy, \textit{Probabilistic Machine Learning: An Introduction}, Ch. 5.%
}

\newcommand{\AMLUnitIIISources}{%
Zhang, Lipton, Li and Smola, \textit{Dive into Deep Learning}, Ch. 12 (Optimization Algorithms).;%
Deisenroth, Faisal and Ong, \textit{Mathematics for Machine Learning}, Ch. 7.;%
Kingma and Ba, ``Adam: A Method for Stochastic Optimization,'' ICLR 2015.;%
Ruder, ``An overview of gradient descent optimization algorithms,'' arXiv:1609.04747, 2016.%
}

\newcommand{\AMLUnitIVSources}{%
Han, Kamber and Pei, \textit{Data Mining: Concepts and Techniques} (3e), Ch. 3.;%
M\"uller and Guido, \textit{Introduction to Machine Learning with Python}, Ch. 3--4.;%
James, Witten, Hastie, Tibshirani and Taylor, \textit{An Introduction to Statistical Learning with Python}, Ch. 5--6, 12.;%
Scikit-learn documentation (preprocessing, model selection), accessed 2026-08-04.%
}

\newcommand{\AMLUnitVSources}{%
James et al., \textit{An Introduction to Statistical Learning with Python}, Ch. 3--6.;%
Bishop, \textit{Pattern Recognition and Machine Learning}, Ch. 3.;%
Deisenroth, Faisal and Ong, \textit{Mathematics for Machine Learning}, Ch. 9.;%
M\"uller and Guido, \textit{Introduction to Machine Learning with Python}, \S2.3.%
}

\newcommand{\AMLUnitVISources}{%
Han, Kamber and Pei, \textit{Data Mining: Concepts and Techniques} (3e), Ch. 8--9.;%
James et al., \textit{An Introduction to Statistical Learning with Python}, Ch. 8--9.;%
Bishop, \textit{Pattern Recognition and Machine Learning}, Ch. 4--5, 7, 14.;%
Zhang et al., \textit{Dive into Deep Learning}, Ch. 5.%
}

\newcommand{\AMLUnitVIISources}{%
Han, Kamber and Pei, \textit{Data Mining: Concepts and Techniques} (3e), Ch. 10.;%
James et al., \textit{An Introduction to Statistical Learning with Python}, Ch. 12.;%
M\"uller and Guido, \textit{Introduction to Machine Learning with Python}, \S3.5.;%
Ester, Kriegel, Sander and Xu, ``A density-based algorithm for discovering clusters (DBSCAN),'' KDD 1996.%
}

\newcommand{\AMLLabSources}{%
M\"uller and Guido, \textit{Introduction to Machine Learning with Python}, companion notebooks.;%
McKinney, \textit{Python for Data Analysis} (3e), O'Reilly, 2022.;%
Scikit-learn user guide, accessed 2026-08-04.;%
Matplotlib and Seaborn documentation, accessed 2026-08-04.%
}
%
    }{%
      % Marking schemes live at `private/solutions/<family>/<id>/latex/`.
      \IfFileExists{../../../../../public/_shared/aml-sources.tex}{%
        % Source strings for Applied Machine Learning (CSAI2017P) — used by slides + notes.
% Prescribed texts and reference books are taken from the course syllabus.
% Support texts are the author-hosted open-access editions:
%   statlearning.com            (ISL with Python)
%   probml.github.io/pml-book   (Murphy, Probabilistic Machine Learning)
%   d2l.ai                      (Dive into Deep Learning)
%   mml-book.github.io          (Mathematics for Machine Learning)

\newcommand{\AMLRepoURL}{https://github.com/tali7c/Applied-Machine-Learning}

% --- prescribed -------------------------------------------------------------
\newcommand{\AMLPrescribed}{%
M\"uller and Guido, \textit{Introduction to Machine Learning with Python}, O'Reilly, 2016.;%
Bishop, \textit{Pattern Recognition and Machine Learning}, Springer, 2006.;%
Murphy, \textit{Machine Learning: A Probabilistic Perspective}, MIT Press, 2012.;%
Han, Kamber and Pei, \textit{Data Mining: Concepts and Techniques} (3e), Morgan Kaufmann, 2012.%
}

% --- unit-wise ---------------------------------------------------------------
\newcommand{\AMLUnitISources}{%
M\"uller and Guido, \textit{Introduction to Machine Learning with Python}, Ch. 1.;%
Bishop, \textit{Pattern Recognition and Machine Learning}, Ch. 1.;%
Murphy, \textit{Probabilistic Machine Learning: An Introduction}, MIT Press, 2022, Ch. 1.;%
Mitchell, \textit{Machine Learning}, McGraw Hill, 1997, Ch. 1.;%
Scikit-learn documentation, accessed 2026-08-04.%
}

\newcommand{\AMLUnitIISources}{%
Bishop, \textit{Pattern Recognition and Machine Learning}, \S1.5, \S4.3, \S7.1.;%
Zhang, Lipton, Li and Smola, \textit{Dive into Deep Learning}, \S3.4.;%
Shalev-Shwartz and Ben-David, \textit{Understanding Machine Learning}, Ch. 15.;%
Murphy, \textit{Probabilistic Machine Learning: An Introduction}, Ch. 5.%
}

\newcommand{\AMLUnitIIISources}{%
Zhang, Lipton, Li and Smola, \textit{Dive into Deep Learning}, Ch. 12 (Optimization Algorithms).;%
Deisenroth, Faisal and Ong, \textit{Mathematics for Machine Learning}, Ch. 7.;%
Kingma and Ba, ``Adam: A Method for Stochastic Optimization,'' ICLR 2015.;%
Ruder, ``An overview of gradient descent optimization algorithms,'' arXiv:1609.04747, 2016.%
}

\newcommand{\AMLUnitIVSources}{%
Han, Kamber and Pei, \textit{Data Mining: Concepts and Techniques} (3e), Ch. 3.;%
M\"uller and Guido, \textit{Introduction to Machine Learning with Python}, Ch. 3--4.;%
James, Witten, Hastie, Tibshirani and Taylor, \textit{An Introduction to Statistical Learning with Python}, Ch. 5--6, 12.;%
Scikit-learn documentation (preprocessing, model selection), accessed 2026-08-04.%
}

\newcommand{\AMLUnitVSources}{%
James et al., \textit{An Introduction to Statistical Learning with Python}, Ch. 3--6.;%
Bishop, \textit{Pattern Recognition and Machine Learning}, Ch. 3.;%
Deisenroth, Faisal and Ong, \textit{Mathematics for Machine Learning}, Ch. 9.;%
M\"uller and Guido, \textit{Introduction to Machine Learning with Python}, \S2.3.%
}

\newcommand{\AMLUnitVISources}{%
Han, Kamber and Pei, \textit{Data Mining: Concepts and Techniques} (3e), Ch. 8--9.;%
James et al., \textit{An Introduction to Statistical Learning with Python}, Ch. 8--9.;%
Bishop, \textit{Pattern Recognition and Machine Learning}, Ch. 4--5, 7, 14.;%
Zhang et al., \textit{Dive into Deep Learning}, Ch. 5.%
}

\newcommand{\AMLUnitVIISources}{%
Han, Kamber and Pei, \textit{Data Mining: Concepts and Techniques} (3e), Ch. 10.;%
James et al., \textit{An Introduction to Statistical Learning with Python}, Ch. 12.;%
M\"uller and Guido, \textit{Introduction to Machine Learning with Python}, \S3.5.;%
Ester, Kriegel, Sander and Xu, ``A density-based algorithm for discovering clusters (DBSCAN),'' KDD 1996.%
}

\newcommand{\AMLLabSources}{%
M\"uller and Guido, \textit{Introduction to Machine Learning with Python}, companion notebooks.;%
McKinney, \textit{Python for Data Analysis} (3e), O'Reilly, 2022.;%
Scikit-learn user guide, accessed 2026-08-04.;%
Matplotlib and Seaborn documentation, accessed 2026-08-04.%
}
%
      }{%
        % Source strings for Applied Machine Learning (CSAI2017P) — used by slides + notes.
% Prescribed texts and reference books are taken from the course syllabus.
% Support texts are the author-hosted open-access editions:
%   statlearning.com            (ISL with Python)
%   probml.github.io/pml-book   (Murphy, Probabilistic Machine Learning)
%   d2l.ai                      (Dive into Deep Learning)
%   mml-book.github.io          (Mathematics for Machine Learning)

\newcommand{\AMLRepoURL}{https://github.com/tali7c/Applied-Machine-Learning}

% --- prescribed -------------------------------------------------------------
\newcommand{\AMLPrescribed}{%
M\"uller and Guido, \textit{Introduction to Machine Learning with Python}, O'Reilly, 2016.;%
Bishop, \textit{Pattern Recognition and Machine Learning}, Springer, 2006.;%
Murphy, \textit{Machine Learning: A Probabilistic Perspective}, MIT Press, 2012.;%
Han, Kamber and Pei, \textit{Data Mining: Concepts and Techniques} (3e), Morgan Kaufmann, 2012.%
}

% --- unit-wise ---------------------------------------------------------------
\newcommand{\AMLUnitISources}{%
M\"uller and Guido, \textit{Introduction to Machine Learning with Python}, Ch. 1.;%
Bishop, \textit{Pattern Recognition and Machine Learning}, Ch. 1.;%
Murphy, \textit{Probabilistic Machine Learning: An Introduction}, MIT Press, 2022, Ch. 1.;%
Mitchell, \textit{Machine Learning}, McGraw Hill, 1997, Ch. 1.;%
Scikit-learn documentation, accessed 2026-08-04.%
}

\newcommand{\AMLUnitIISources}{%
Bishop, \textit{Pattern Recognition and Machine Learning}, \S1.5, \S4.3, \S7.1.;%
Zhang, Lipton, Li and Smola, \textit{Dive into Deep Learning}, \S3.4.;%
Shalev-Shwartz and Ben-David, \textit{Understanding Machine Learning}, Ch. 15.;%
Murphy, \textit{Probabilistic Machine Learning: An Introduction}, Ch. 5.%
}

\newcommand{\AMLUnitIIISources}{%
Zhang, Lipton, Li and Smola, \textit{Dive into Deep Learning}, Ch. 12 (Optimization Algorithms).;%
Deisenroth, Faisal and Ong, \textit{Mathematics for Machine Learning}, Ch. 7.;%
Kingma and Ba, ``Adam: A Method for Stochastic Optimization,'' ICLR 2015.;%
Ruder, ``An overview of gradient descent optimization algorithms,'' arXiv:1609.04747, 2016.%
}

\newcommand{\AMLUnitIVSources}{%
Han, Kamber and Pei, \textit{Data Mining: Concepts and Techniques} (3e), Ch. 3.;%
M\"uller and Guido, \textit{Introduction to Machine Learning with Python}, Ch. 3--4.;%
James, Witten, Hastie, Tibshirani and Taylor, \textit{An Introduction to Statistical Learning with Python}, Ch. 5--6, 12.;%
Scikit-learn documentation (preprocessing, model selection), accessed 2026-08-04.%
}

\newcommand{\AMLUnitVSources}{%
James et al., \textit{An Introduction to Statistical Learning with Python}, Ch. 3--6.;%
Bishop, \textit{Pattern Recognition and Machine Learning}, Ch. 3.;%
Deisenroth, Faisal and Ong, \textit{Mathematics for Machine Learning}, Ch. 9.;%
M\"uller and Guido, \textit{Introduction to Machine Learning with Python}, \S2.3.%
}

\newcommand{\AMLUnitVISources}{%
Han, Kamber and Pei, \textit{Data Mining: Concepts and Techniques} (3e), Ch. 8--9.;%
James et al., \textit{An Introduction to Statistical Learning with Python}, Ch. 8--9.;%
Bishop, \textit{Pattern Recognition and Machine Learning}, Ch. 4--5, 7, 14.;%
Zhang et al., \textit{Dive into Deep Learning}, Ch. 5.%
}

\newcommand{\AMLUnitVIISources}{%
Han, Kamber and Pei, \textit{Data Mining: Concepts and Techniques} (3e), Ch. 10.;%
James et al., \textit{An Introduction to Statistical Learning with Python}, Ch. 12.;%
M\"uller and Guido, \textit{Introduction to Machine Learning with Python}, \S3.5.;%
Ester, Kriegel, Sander and Xu, ``A density-based algorithm for discovering clusters (DBSCAN),'' KDD 1996.%
}

\newcommand{\AMLLabSources}{%
M\"uller and Guido, \textit{Introduction to Machine Learning with Python}, companion notebooks.;%
McKinney, \textit{Python for Data Analysis} (3e), O'Reilly, 2022.;%
Scikit-learn user guide, accessed 2026-08-04.;%
Matplotlib and Seaborn documentation, accessed 2026-08-04.%
}
%
      }%
    }%
  }%
}
%
}{%
  \IfFileExists{../../../../../public/_shared/notes-common.tex}{%
    % Shared notes settings for Applied Machine Learning lectures (UPES).
% Keep this file free of \documentclass to allow reuse via \input.

\usepackage[utf8]{inputenc}
\usepackage[T1]{fontenc}
\usepackage{lmodern}          % scalable T1 fonts; without this pdflatex falls back
\input{glyphtounicode}        % to bitmapped EC Type 3 and the PDFs are neither
\pdfgentounicode=1            % searchable nor screen-reader friendly
\usepackage[a4paper,margin=1in]{geometry}
\usepackage{hyperref}
\usepackage{xurl}
\usepackage{booktabs}
\usepackage{amsmath}
\usepackage{amssymb}
\usepackage{listings}
\usepackage{xcolor}
\usepackage{enumitem}
\usepackage{parskip}

% Code listings (Python-oriented for this subject).
\lstset{
  language=Python,
  basicstyle=\ttfamily\small,
  keywordstyle=\color{blue},
  commentstyle=\color{gray},
  stringstyle=\color{teal},
  showstringspaces=false,
  columns=fullflexible,
  frame=single,
  framerule=0.2pt,
  breaklines=true
}

\setlist[itemize]{topsep=0.3em,itemsep=0.2em}
\setlist[enumerate]{topsep=0.3em,itemsep=0.2em}

% Simple per-section source line.
\newcommand{\notesources}[1]{\par\noindent{\small\color{gray}\textit{Sources:} \sloppy #1\par}}

% Unit-wise source strings for this subject (used by slides + notes).
% Path is relative to the lecture `latex/` folder (the compilation CWD).
\IfFileExists{../../../_shared/aml-sources.tex}{%
  % Source strings for Applied Machine Learning (CSAI2017P) — used by slides + notes.
% Prescribed texts and reference books are taken from the course syllabus.
% Support texts are the author-hosted open-access editions:
%   statlearning.com            (ISL with Python)
%   probml.github.io/pml-book   (Murphy, Probabilistic Machine Learning)
%   d2l.ai                      (Dive into Deep Learning)
%   mml-book.github.io          (Mathematics for Machine Learning)

\newcommand{\AMLRepoURL}{https://github.com/tali7c/Applied-Machine-Learning}

% --- prescribed -------------------------------------------------------------
\newcommand{\AMLPrescribed}{%
M\"uller and Guido, \textit{Introduction to Machine Learning with Python}, O'Reilly, 2016.;%
Bishop, \textit{Pattern Recognition and Machine Learning}, Springer, 2006.;%
Murphy, \textit{Machine Learning: A Probabilistic Perspective}, MIT Press, 2012.;%
Han, Kamber and Pei, \textit{Data Mining: Concepts and Techniques} (3e), Morgan Kaufmann, 2012.%
}

% --- unit-wise ---------------------------------------------------------------
\newcommand{\AMLUnitISources}{%
M\"uller and Guido, \textit{Introduction to Machine Learning with Python}, Ch. 1.;%
Bishop, \textit{Pattern Recognition and Machine Learning}, Ch. 1.;%
Murphy, \textit{Probabilistic Machine Learning: An Introduction}, MIT Press, 2022, Ch. 1.;%
Mitchell, \textit{Machine Learning}, McGraw Hill, 1997, Ch. 1.;%
Scikit-learn documentation, accessed 2026-08-04.%
}

\newcommand{\AMLUnitIISources}{%
Bishop, \textit{Pattern Recognition and Machine Learning}, \S1.5, \S4.3, \S7.1.;%
Zhang, Lipton, Li and Smola, \textit{Dive into Deep Learning}, \S3.4.;%
Shalev-Shwartz and Ben-David, \textit{Understanding Machine Learning}, Ch. 15.;%
Murphy, \textit{Probabilistic Machine Learning: An Introduction}, Ch. 5.%
}

\newcommand{\AMLUnitIIISources}{%
Zhang, Lipton, Li and Smola, \textit{Dive into Deep Learning}, Ch. 12 (Optimization Algorithms).;%
Deisenroth, Faisal and Ong, \textit{Mathematics for Machine Learning}, Ch. 7.;%
Kingma and Ba, ``Adam: A Method for Stochastic Optimization,'' ICLR 2015.;%
Ruder, ``An overview of gradient descent optimization algorithms,'' arXiv:1609.04747, 2016.%
}

\newcommand{\AMLUnitIVSources}{%
Han, Kamber and Pei, \textit{Data Mining: Concepts and Techniques} (3e), Ch. 3.;%
M\"uller and Guido, \textit{Introduction to Machine Learning with Python}, Ch. 3--4.;%
James, Witten, Hastie, Tibshirani and Taylor, \textit{An Introduction to Statistical Learning with Python}, Ch. 5--6, 12.;%
Scikit-learn documentation (preprocessing, model selection), accessed 2026-08-04.%
}

\newcommand{\AMLUnitVSources}{%
James et al., \textit{An Introduction to Statistical Learning with Python}, Ch. 3--6.;%
Bishop, \textit{Pattern Recognition and Machine Learning}, Ch. 3.;%
Deisenroth, Faisal and Ong, \textit{Mathematics for Machine Learning}, Ch. 9.;%
M\"uller and Guido, \textit{Introduction to Machine Learning with Python}, \S2.3.%
}

\newcommand{\AMLUnitVISources}{%
Han, Kamber and Pei, \textit{Data Mining: Concepts and Techniques} (3e), Ch. 8--9.;%
James et al., \textit{An Introduction to Statistical Learning with Python}, Ch. 8--9.;%
Bishop, \textit{Pattern Recognition and Machine Learning}, Ch. 4--5, 7, 14.;%
Zhang et al., \textit{Dive into Deep Learning}, Ch. 5.%
}

\newcommand{\AMLUnitVIISources}{%
Han, Kamber and Pei, \textit{Data Mining: Concepts and Techniques} (3e), Ch. 10.;%
James et al., \textit{An Introduction to Statistical Learning with Python}, Ch. 12.;%
M\"uller and Guido, \textit{Introduction to Machine Learning with Python}, \S3.5.;%
Ester, Kriegel, Sander and Xu, ``A density-based algorithm for discovering clusters (DBSCAN),'' KDD 1996.%
}

\newcommand{\AMLLabSources}{%
M\"uller and Guido, \textit{Introduction to Machine Learning with Python}, companion notebooks.;%
McKinney, \textit{Python for Data Analysis} (3e), O'Reilly, 2022.;%
Scikit-learn user guide, accessed 2026-08-04.;%
Matplotlib and Seaborn documentation, accessed 2026-08-04.%
}
%
}{%
  % Faculty copies live under `private/` so their compile CWD is different.
  \IfFileExists{../../../../public/_shared/aml-sources.tex}{%
    % Source strings for Applied Machine Learning (CSAI2017P) — used by slides + notes.
% Prescribed texts and reference books are taken from the course syllabus.
% Support texts are the author-hosted open-access editions:
%   statlearning.com            (ISL with Python)
%   probml.github.io/pml-book   (Murphy, Probabilistic Machine Learning)
%   d2l.ai                      (Dive into Deep Learning)
%   mml-book.github.io          (Mathematics for Machine Learning)

\newcommand{\AMLRepoURL}{https://github.com/tali7c/Applied-Machine-Learning}

% --- prescribed -------------------------------------------------------------
\newcommand{\AMLPrescribed}{%
M\"uller and Guido, \textit{Introduction to Machine Learning with Python}, O'Reilly, 2016.;%
Bishop, \textit{Pattern Recognition and Machine Learning}, Springer, 2006.;%
Murphy, \textit{Machine Learning: A Probabilistic Perspective}, MIT Press, 2012.;%
Han, Kamber and Pei, \textit{Data Mining: Concepts and Techniques} (3e), Morgan Kaufmann, 2012.%
}

% --- unit-wise ---------------------------------------------------------------
\newcommand{\AMLUnitISources}{%
M\"uller and Guido, \textit{Introduction to Machine Learning with Python}, Ch. 1.;%
Bishop, \textit{Pattern Recognition and Machine Learning}, Ch. 1.;%
Murphy, \textit{Probabilistic Machine Learning: An Introduction}, MIT Press, 2022, Ch. 1.;%
Mitchell, \textit{Machine Learning}, McGraw Hill, 1997, Ch. 1.;%
Scikit-learn documentation, accessed 2026-08-04.%
}

\newcommand{\AMLUnitIISources}{%
Bishop, \textit{Pattern Recognition and Machine Learning}, \S1.5, \S4.3, \S7.1.;%
Zhang, Lipton, Li and Smola, \textit{Dive into Deep Learning}, \S3.4.;%
Shalev-Shwartz and Ben-David, \textit{Understanding Machine Learning}, Ch. 15.;%
Murphy, \textit{Probabilistic Machine Learning: An Introduction}, Ch. 5.%
}

\newcommand{\AMLUnitIIISources}{%
Zhang, Lipton, Li and Smola, \textit{Dive into Deep Learning}, Ch. 12 (Optimization Algorithms).;%
Deisenroth, Faisal and Ong, \textit{Mathematics for Machine Learning}, Ch. 7.;%
Kingma and Ba, ``Adam: A Method for Stochastic Optimization,'' ICLR 2015.;%
Ruder, ``An overview of gradient descent optimization algorithms,'' arXiv:1609.04747, 2016.%
}

\newcommand{\AMLUnitIVSources}{%
Han, Kamber and Pei, \textit{Data Mining: Concepts and Techniques} (3e), Ch. 3.;%
M\"uller and Guido, \textit{Introduction to Machine Learning with Python}, Ch. 3--4.;%
James, Witten, Hastie, Tibshirani and Taylor, \textit{An Introduction to Statistical Learning with Python}, Ch. 5--6, 12.;%
Scikit-learn documentation (preprocessing, model selection), accessed 2026-08-04.%
}

\newcommand{\AMLUnitVSources}{%
James et al., \textit{An Introduction to Statistical Learning with Python}, Ch. 3--6.;%
Bishop, \textit{Pattern Recognition and Machine Learning}, Ch. 3.;%
Deisenroth, Faisal and Ong, \textit{Mathematics for Machine Learning}, Ch. 9.;%
M\"uller and Guido, \textit{Introduction to Machine Learning with Python}, \S2.3.%
}

\newcommand{\AMLUnitVISources}{%
Han, Kamber and Pei, \textit{Data Mining: Concepts and Techniques} (3e), Ch. 8--9.;%
James et al., \textit{An Introduction to Statistical Learning with Python}, Ch. 8--9.;%
Bishop, \textit{Pattern Recognition and Machine Learning}, Ch. 4--5, 7, 14.;%
Zhang et al., \textit{Dive into Deep Learning}, Ch. 5.%
}

\newcommand{\AMLUnitVIISources}{%
Han, Kamber and Pei, \textit{Data Mining: Concepts and Techniques} (3e), Ch. 10.;%
James et al., \textit{An Introduction to Statistical Learning with Python}, Ch. 12.;%
M\"uller and Guido, \textit{Introduction to Machine Learning with Python}, \S3.5.;%
Ester, Kriegel, Sander and Xu, ``A density-based algorithm for discovering clusters (DBSCAN),'' KDD 1996.%
}

\newcommand{\AMLLabSources}{%
M\"uller and Guido, \textit{Introduction to Machine Learning with Python}, companion notebooks.;%
McKinney, \textit{Python for Data Analysis} (3e), O'Reilly, 2022.;%
Scikit-learn user guide, accessed 2026-08-04.;%
Matplotlib and Seaborn documentation, accessed 2026-08-04.%
}
%
  }{%
    % Question papers live at `private/<family>/<instrument>/`.
    \IfFileExists{../../../public/_shared/aml-sources.tex}{%
      % Source strings for Applied Machine Learning (CSAI2017P) — used by slides + notes.
% Prescribed texts and reference books are taken from the course syllabus.
% Support texts are the author-hosted open-access editions:
%   statlearning.com            (ISL with Python)
%   probml.github.io/pml-book   (Murphy, Probabilistic Machine Learning)
%   d2l.ai                      (Dive into Deep Learning)
%   mml-book.github.io          (Mathematics for Machine Learning)

\newcommand{\AMLRepoURL}{https://github.com/tali7c/Applied-Machine-Learning}

% --- prescribed -------------------------------------------------------------
\newcommand{\AMLPrescribed}{%
M\"uller and Guido, \textit{Introduction to Machine Learning with Python}, O'Reilly, 2016.;%
Bishop, \textit{Pattern Recognition and Machine Learning}, Springer, 2006.;%
Murphy, \textit{Machine Learning: A Probabilistic Perspective}, MIT Press, 2012.;%
Han, Kamber and Pei, \textit{Data Mining: Concepts and Techniques} (3e), Morgan Kaufmann, 2012.%
}

% --- unit-wise ---------------------------------------------------------------
\newcommand{\AMLUnitISources}{%
M\"uller and Guido, \textit{Introduction to Machine Learning with Python}, Ch. 1.;%
Bishop, \textit{Pattern Recognition and Machine Learning}, Ch. 1.;%
Murphy, \textit{Probabilistic Machine Learning: An Introduction}, MIT Press, 2022, Ch. 1.;%
Mitchell, \textit{Machine Learning}, McGraw Hill, 1997, Ch. 1.;%
Scikit-learn documentation, accessed 2026-08-04.%
}

\newcommand{\AMLUnitIISources}{%
Bishop, \textit{Pattern Recognition and Machine Learning}, \S1.5, \S4.3, \S7.1.;%
Zhang, Lipton, Li and Smola, \textit{Dive into Deep Learning}, \S3.4.;%
Shalev-Shwartz and Ben-David, \textit{Understanding Machine Learning}, Ch. 15.;%
Murphy, \textit{Probabilistic Machine Learning: An Introduction}, Ch. 5.%
}

\newcommand{\AMLUnitIIISources}{%
Zhang, Lipton, Li and Smola, \textit{Dive into Deep Learning}, Ch. 12 (Optimization Algorithms).;%
Deisenroth, Faisal and Ong, \textit{Mathematics for Machine Learning}, Ch. 7.;%
Kingma and Ba, ``Adam: A Method for Stochastic Optimization,'' ICLR 2015.;%
Ruder, ``An overview of gradient descent optimization algorithms,'' arXiv:1609.04747, 2016.%
}

\newcommand{\AMLUnitIVSources}{%
Han, Kamber and Pei, \textit{Data Mining: Concepts and Techniques} (3e), Ch. 3.;%
M\"uller and Guido, \textit{Introduction to Machine Learning with Python}, Ch. 3--4.;%
James, Witten, Hastie, Tibshirani and Taylor, \textit{An Introduction to Statistical Learning with Python}, Ch. 5--6, 12.;%
Scikit-learn documentation (preprocessing, model selection), accessed 2026-08-04.%
}

\newcommand{\AMLUnitVSources}{%
James et al., \textit{An Introduction to Statistical Learning with Python}, Ch. 3--6.;%
Bishop, \textit{Pattern Recognition and Machine Learning}, Ch. 3.;%
Deisenroth, Faisal and Ong, \textit{Mathematics for Machine Learning}, Ch. 9.;%
M\"uller and Guido, \textit{Introduction to Machine Learning with Python}, \S2.3.%
}

\newcommand{\AMLUnitVISources}{%
Han, Kamber and Pei, \textit{Data Mining: Concepts and Techniques} (3e), Ch. 8--9.;%
James et al., \textit{An Introduction to Statistical Learning with Python}, Ch. 8--9.;%
Bishop, \textit{Pattern Recognition and Machine Learning}, Ch. 4--5, 7, 14.;%
Zhang et al., \textit{Dive into Deep Learning}, Ch. 5.%
}

\newcommand{\AMLUnitVIISources}{%
Han, Kamber and Pei, \textit{Data Mining: Concepts and Techniques} (3e), Ch. 10.;%
James et al., \textit{An Introduction to Statistical Learning with Python}, Ch. 12.;%
M\"uller and Guido, \textit{Introduction to Machine Learning with Python}, \S3.5.;%
Ester, Kriegel, Sander and Xu, ``A density-based algorithm for discovering clusters (DBSCAN),'' KDD 1996.%
}

\newcommand{\AMLLabSources}{%
M\"uller and Guido, \textit{Introduction to Machine Learning with Python}, companion notebooks.;%
McKinney, \textit{Python for Data Analysis} (3e), O'Reilly, 2022.;%
Scikit-learn user guide, accessed 2026-08-04.;%
Matplotlib and Seaborn documentation, accessed 2026-08-04.%
}
%
    }{%
      % Marking schemes live at `private/solutions/<family>/<id>/latex/`.
      \IfFileExists{../../../../../public/_shared/aml-sources.tex}{%
        % Source strings for Applied Machine Learning (CSAI2017P) — used by slides + notes.
% Prescribed texts and reference books are taken from the course syllabus.
% Support texts are the author-hosted open-access editions:
%   statlearning.com            (ISL with Python)
%   probml.github.io/pml-book   (Murphy, Probabilistic Machine Learning)
%   d2l.ai                      (Dive into Deep Learning)
%   mml-book.github.io          (Mathematics for Machine Learning)

\newcommand{\AMLRepoURL}{https://github.com/tali7c/Applied-Machine-Learning}

% --- prescribed -------------------------------------------------------------
\newcommand{\AMLPrescribed}{%
M\"uller and Guido, \textit{Introduction to Machine Learning with Python}, O'Reilly, 2016.;%
Bishop, \textit{Pattern Recognition and Machine Learning}, Springer, 2006.;%
Murphy, \textit{Machine Learning: A Probabilistic Perspective}, MIT Press, 2012.;%
Han, Kamber and Pei, \textit{Data Mining: Concepts and Techniques} (3e), Morgan Kaufmann, 2012.%
}

% --- unit-wise ---------------------------------------------------------------
\newcommand{\AMLUnitISources}{%
M\"uller and Guido, \textit{Introduction to Machine Learning with Python}, Ch. 1.;%
Bishop, \textit{Pattern Recognition and Machine Learning}, Ch. 1.;%
Murphy, \textit{Probabilistic Machine Learning: An Introduction}, MIT Press, 2022, Ch. 1.;%
Mitchell, \textit{Machine Learning}, McGraw Hill, 1997, Ch. 1.;%
Scikit-learn documentation, accessed 2026-08-04.%
}

\newcommand{\AMLUnitIISources}{%
Bishop, \textit{Pattern Recognition and Machine Learning}, \S1.5, \S4.3, \S7.1.;%
Zhang, Lipton, Li and Smola, \textit{Dive into Deep Learning}, \S3.4.;%
Shalev-Shwartz and Ben-David, \textit{Understanding Machine Learning}, Ch. 15.;%
Murphy, \textit{Probabilistic Machine Learning: An Introduction}, Ch. 5.%
}

\newcommand{\AMLUnitIIISources}{%
Zhang, Lipton, Li and Smola, \textit{Dive into Deep Learning}, Ch. 12 (Optimization Algorithms).;%
Deisenroth, Faisal and Ong, \textit{Mathematics for Machine Learning}, Ch. 7.;%
Kingma and Ba, ``Adam: A Method for Stochastic Optimization,'' ICLR 2015.;%
Ruder, ``An overview of gradient descent optimization algorithms,'' arXiv:1609.04747, 2016.%
}

\newcommand{\AMLUnitIVSources}{%
Han, Kamber and Pei, \textit{Data Mining: Concepts and Techniques} (3e), Ch. 3.;%
M\"uller and Guido, \textit{Introduction to Machine Learning with Python}, Ch. 3--4.;%
James, Witten, Hastie, Tibshirani and Taylor, \textit{An Introduction to Statistical Learning with Python}, Ch. 5--6, 12.;%
Scikit-learn documentation (preprocessing, model selection), accessed 2026-08-04.%
}

\newcommand{\AMLUnitVSources}{%
James et al., \textit{An Introduction to Statistical Learning with Python}, Ch. 3--6.;%
Bishop, \textit{Pattern Recognition and Machine Learning}, Ch. 3.;%
Deisenroth, Faisal and Ong, \textit{Mathematics for Machine Learning}, Ch. 9.;%
M\"uller and Guido, \textit{Introduction to Machine Learning with Python}, \S2.3.%
}

\newcommand{\AMLUnitVISources}{%
Han, Kamber and Pei, \textit{Data Mining: Concepts and Techniques} (3e), Ch. 8--9.;%
James et al., \textit{An Introduction to Statistical Learning with Python}, Ch. 8--9.;%
Bishop, \textit{Pattern Recognition and Machine Learning}, Ch. 4--5, 7, 14.;%
Zhang et al., \textit{Dive into Deep Learning}, Ch. 5.%
}

\newcommand{\AMLUnitVIISources}{%
Han, Kamber and Pei, \textit{Data Mining: Concepts and Techniques} (3e), Ch. 10.;%
James et al., \textit{An Introduction to Statistical Learning with Python}, Ch. 12.;%
M\"uller and Guido, \textit{Introduction to Machine Learning with Python}, \S3.5.;%
Ester, Kriegel, Sander and Xu, ``A density-based algorithm for discovering clusters (DBSCAN),'' KDD 1996.%
}

\newcommand{\AMLLabSources}{%
M\"uller and Guido, \textit{Introduction to Machine Learning with Python}, companion notebooks.;%
McKinney, \textit{Python for Data Analysis} (3e), O'Reilly, 2022.;%
Scikit-learn user guide, accessed 2026-08-04.;%
Matplotlib and Seaborn documentation, accessed 2026-08-04.%
}
%
      }{%
        % Source strings for Applied Machine Learning (CSAI2017P) — used by slides + notes.
% Prescribed texts and reference books are taken from the course syllabus.
% Support texts are the author-hosted open-access editions:
%   statlearning.com            (ISL with Python)
%   probml.github.io/pml-book   (Murphy, Probabilistic Machine Learning)
%   d2l.ai                      (Dive into Deep Learning)
%   mml-book.github.io          (Mathematics for Machine Learning)

\newcommand{\AMLRepoURL}{https://github.com/tali7c/Applied-Machine-Learning}

% --- prescribed -------------------------------------------------------------
\newcommand{\AMLPrescribed}{%
M\"uller and Guido, \textit{Introduction to Machine Learning with Python}, O'Reilly, 2016.;%
Bishop, \textit{Pattern Recognition and Machine Learning}, Springer, 2006.;%
Murphy, \textit{Machine Learning: A Probabilistic Perspective}, MIT Press, 2012.;%
Han, Kamber and Pei, \textit{Data Mining: Concepts and Techniques} (3e), Morgan Kaufmann, 2012.%
}

% --- unit-wise ---------------------------------------------------------------
\newcommand{\AMLUnitISources}{%
M\"uller and Guido, \textit{Introduction to Machine Learning with Python}, Ch. 1.;%
Bishop, \textit{Pattern Recognition and Machine Learning}, Ch. 1.;%
Murphy, \textit{Probabilistic Machine Learning: An Introduction}, MIT Press, 2022, Ch. 1.;%
Mitchell, \textit{Machine Learning}, McGraw Hill, 1997, Ch. 1.;%
Scikit-learn documentation, accessed 2026-08-04.%
}

\newcommand{\AMLUnitIISources}{%
Bishop, \textit{Pattern Recognition and Machine Learning}, \S1.5, \S4.3, \S7.1.;%
Zhang, Lipton, Li and Smola, \textit{Dive into Deep Learning}, \S3.4.;%
Shalev-Shwartz and Ben-David, \textit{Understanding Machine Learning}, Ch. 15.;%
Murphy, \textit{Probabilistic Machine Learning: An Introduction}, Ch. 5.%
}

\newcommand{\AMLUnitIIISources}{%
Zhang, Lipton, Li and Smola, \textit{Dive into Deep Learning}, Ch. 12 (Optimization Algorithms).;%
Deisenroth, Faisal and Ong, \textit{Mathematics for Machine Learning}, Ch. 7.;%
Kingma and Ba, ``Adam: A Method for Stochastic Optimization,'' ICLR 2015.;%
Ruder, ``An overview of gradient descent optimization algorithms,'' arXiv:1609.04747, 2016.%
}

\newcommand{\AMLUnitIVSources}{%
Han, Kamber and Pei, \textit{Data Mining: Concepts and Techniques} (3e), Ch. 3.;%
M\"uller and Guido, \textit{Introduction to Machine Learning with Python}, Ch. 3--4.;%
James, Witten, Hastie, Tibshirani and Taylor, \textit{An Introduction to Statistical Learning with Python}, Ch. 5--6, 12.;%
Scikit-learn documentation (preprocessing, model selection), accessed 2026-08-04.%
}

\newcommand{\AMLUnitVSources}{%
James et al., \textit{An Introduction to Statistical Learning with Python}, Ch. 3--6.;%
Bishop, \textit{Pattern Recognition and Machine Learning}, Ch. 3.;%
Deisenroth, Faisal and Ong, \textit{Mathematics for Machine Learning}, Ch. 9.;%
M\"uller and Guido, \textit{Introduction to Machine Learning with Python}, \S2.3.%
}

\newcommand{\AMLUnitVISources}{%
Han, Kamber and Pei, \textit{Data Mining: Concepts and Techniques} (3e), Ch. 8--9.;%
James et al., \textit{An Introduction to Statistical Learning with Python}, Ch. 8--9.;%
Bishop, \textit{Pattern Recognition and Machine Learning}, Ch. 4--5, 7, 14.;%
Zhang et al., \textit{Dive into Deep Learning}, Ch. 5.%
}

\newcommand{\AMLUnitVIISources}{%
Han, Kamber and Pei, \textit{Data Mining: Concepts and Techniques} (3e), Ch. 10.;%
James et al., \textit{An Introduction to Statistical Learning with Python}, Ch. 12.;%
M\"uller and Guido, \textit{Introduction to Machine Learning with Python}, \S3.5.;%
Ester, Kriegel, Sander and Xu, ``A density-based algorithm for discovering clusters (DBSCAN),'' KDD 1996.%
}

\newcommand{\AMLLabSources}{%
M\"uller and Guido, \textit{Introduction to Machine Learning with Python}, companion notebooks.;%
McKinney, \textit{Python for Data Analysis} (3e), O'Reilly, 2022.;%
Scikit-learn user guide, accessed 2026-08-04.;%
Matplotlib and Seaborn documentation, accessed 2026-08-04.%
}
%
      }%
    }%
  }%
}
%
  }{%
    % Shared notes settings for Applied Machine Learning lectures (UPES).
% Keep this file free of \documentclass to allow reuse via \input.

\usepackage[utf8]{inputenc}
\usepackage[T1]{fontenc}
\usepackage{lmodern}          % scalable T1 fonts; without this pdflatex falls back
\input{glyphtounicode}        % to bitmapped EC Type 3 and the PDFs are neither
\pdfgentounicode=1            % searchable nor screen-reader friendly
\usepackage[a4paper,margin=1in]{geometry}
\usepackage{hyperref}
\usepackage{xurl}
\usepackage{booktabs}
\usepackage{amsmath}
\usepackage{amssymb}
\usepackage{listings}
\usepackage{xcolor}
\usepackage{enumitem}
\usepackage{parskip}

% Code listings (Python-oriented for this subject).
\lstset{
  language=Python,
  basicstyle=\ttfamily\small,
  keywordstyle=\color{blue},
  commentstyle=\color{gray},
  stringstyle=\color{teal},
  showstringspaces=false,
  columns=fullflexible,
  frame=single,
  framerule=0.2pt,
  breaklines=true
}

\setlist[itemize]{topsep=0.3em,itemsep=0.2em}
\setlist[enumerate]{topsep=0.3em,itemsep=0.2em}

% Simple per-section source line.
\newcommand{\notesources}[1]{\par\noindent{\small\color{gray}\textit{Sources:} \sloppy #1\par}}

% Unit-wise source strings for this subject (used by slides + notes).
% Path is relative to the lecture `latex/` folder (the compilation CWD).
\IfFileExists{../../../_shared/aml-sources.tex}{%
  % Source strings for Applied Machine Learning (CSAI2017P) — used by slides + notes.
% Prescribed texts and reference books are taken from the course syllabus.
% Support texts are the author-hosted open-access editions:
%   statlearning.com            (ISL with Python)
%   probml.github.io/pml-book   (Murphy, Probabilistic Machine Learning)
%   d2l.ai                      (Dive into Deep Learning)
%   mml-book.github.io          (Mathematics for Machine Learning)

\newcommand{\AMLRepoURL}{https://github.com/tali7c/Applied-Machine-Learning}

% --- prescribed -------------------------------------------------------------
\newcommand{\AMLPrescribed}{%
M\"uller and Guido, \textit{Introduction to Machine Learning with Python}, O'Reilly, 2016.;%
Bishop, \textit{Pattern Recognition and Machine Learning}, Springer, 2006.;%
Murphy, \textit{Machine Learning: A Probabilistic Perspective}, MIT Press, 2012.;%
Han, Kamber and Pei, \textit{Data Mining: Concepts and Techniques} (3e), Morgan Kaufmann, 2012.%
}

% --- unit-wise ---------------------------------------------------------------
\newcommand{\AMLUnitISources}{%
M\"uller and Guido, \textit{Introduction to Machine Learning with Python}, Ch. 1.;%
Bishop, \textit{Pattern Recognition and Machine Learning}, Ch. 1.;%
Murphy, \textit{Probabilistic Machine Learning: An Introduction}, MIT Press, 2022, Ch. 1.;%
Mitchell, \textit{Machine Learning}, McGraw Hill, 1997, Ch. 1.;%
Scikit-learn documentation, accessed 2026-08-04.%
}

\newcommand{\AMLUnitIISources}{%
Bishop, \textit{Pattern Recognition and Machine Learning}, \S1.5, \S4.3, \S7.1.;%
Zhang, Lipton, Li and Smola, \textit{Dive into Deep Learning}, \S3.4.;%
Shalev-Shwartz and Ben-David, \textit{Understanding Machine Learning}, Ch. 15.;%
Murphy, \textit{Probabilistic Machine Learning: An Introduction}, Ch. 5.%
}

\newcommand{\AMLUnitIIISources}{%
Zhang, Lipton, Li and Smola, \textit{Dive into Deep Learning}, Ch. 12 (Optimization Algorithms).;%
Deisenroth, Faisal and Ong, \textit{Mathematics for Machine Learning}, Ch. 7.;%
Kingma and Ba, ``Adam: A Method for Stochastic Optimization,'' ICLR 2015.;%
Ruder, ``An overview of gradient descent optimization algorithms,'' arXiv:1609.04747, 2016.%
}

\newcommand{\AMLUnitIVSources}{%
Han, Kamber and Pei, \textit{Data Mining: Concepts and Techniques} (3e), Ch. 3.;%
M\"uller and Guido, \textit{Introduction to Machine Learning with Python}, Ch. 3--4.;%
James, Witten, Hastie, Tibshirani and Taylor, \textit{An Introduction to Statistical Learning with Python}, Ch. 5--6, 12.;%
Scikit-learn documentation (preprocessing, model selection), accessed 2026-08-04.%
}

\newcommand{\AMLUnitVSources}{%
James et al., \textit{An Introduction to Statistical Learning with Python}, Ch. 3--6.;%
Bishop, \textit{Pattern Recognition and Machine Learning}, Ch. 3.;%
Deisenroth, Faisal and Ong, \textit{Mathematics for Machine Learning}, Ch. 9.;%
M\"uller and Guido, \textit{Introduction to Machine Learning with Python}, \S2.3.%
}

\newcommand{\AMLUnitVISources}{%
Han, Kamber and Pei, \textit{Data Mining: Concepts and Techniques} (3e), Ch. 8--9.;%
James et al., \textit{An Introduction to Statistical Learning with Python}, Ch. 8--9.;%
Bishop, \textit{Pattern Recognition and Machine Learning}, Ch. 4--5, 7, 14.;%
Zhang et al., \textit{Dive into Deep Learning}, Ch. 5.%
}

\newcommand{\AMLUnitVIISources}{%
Han, Kamber and Pei, \textit{Data Mining: Concepts and Techniques} (3e), Ch. 10.;%
James et al., \textit{An Introduction to Statistical Learning with Python}, Ch. 12.;%
M\"uller and Guido, \textit{Introduction to Machine Learning with Python}, \S3.5.;%
Ester, Kriegel, Sander and Xu, ``A density-based algorithm for discovering clusters (DBSCAN),'' KDD 1996.%
}

\newcommand{\AMLLabSources}{%
M\"uller and Guido, \textit{Introduction to Machine Learning with Python}, companion notebooks.;%
McKinney, \textit{Python for Data Analysis} (3e), O'Reilly, 2022.;%
Scikit-learn user guide, accessed 2026-08-04.;%
Matplotlib and Seaborn documentation, accessed 2026-08-04.%
}
%
}{%
  % Faculty copies live under `private/` so their compile CWD is different.
  \IfFileExists{../../../../public/_shared/aml-sources.tex}{%
    % Source strings for Applied Machine Learning (CSAI2017P) — used by slides + notes.
% Prescribed texts and reference books are taken from the course syllabus.
% Support texts are the author-hosted open-access editions:
%   statlearning.com            (ISL with Python)
%   probml.github.io/pml-book   (Murphy, Probabilistic Machine Learning)
%   d2l.ai                      (Dive into Deep Learning)
%   mml-book.github.io          (Mathematics for Machine Learning)

\newcommand{\AMLRepoURL}{https://github.com/tali7c/Applied-Machine-Learning}

% --- prescribed -------------------------------------------------------------
\newcommand{\AMLPrescribed}{%
M\"uller and Guido, \textit{Introduction to Machine Learning with Python}, O'Reilly, 2016.;%
Bishop, \textit{Pattern Recognition and Machine Learning}, Springer, 2006.;%
Murphy, \textit{Machine Learning: A Probabilistic Perspective}, MIT Press, 2012.;%
Han, Kamber and Pei, \textit{Data Mining: Concepts and Techniques} (3e), Morgan Kaufmann, 2012.%
}

% --- unit-wise ---------------------------------------------------------------
\newcommand{\AMLUnitISources}{%
M\"uller and Guido, \textit{Introduction to Machine Learning with Python}, Ch. 1.;%
Bishop, \textit{Pattern Recognition and Machine Learning}, Ch. 1.;%
Murphy, \textit{Probabilistic Machine Learning: An Introduction}, MIT Press, 2022, Ch. 1.;%
Mitchell, \textit{Machine Learning}, McGraw Hill, 1997, Ch. 1.;%
Scikit-learn documentation, accessed 2026-08-04.%
}

\newcommand{\AMLUnitIISources}{%
Bishop, \textit{Pattern Recognition and Machine Learning}, \S1.5, \S4.3, \S7.1.;%
Zhang, Lipton, Li and Smola, \textit{Dive into Deep Learning}, \S3.4.;%
Shalev-Shwartz and Ben-David, \textit{Understanding Machine Learning}, Ch. 15.;%
Murphy, \textit{Probabilistic Machine Learning: An Introduction}, Ch. 5.%
}

\newcommand{\AMLUnitIIISources}{%
Zhang, Lipton, Li and Smola, \textit{Dive into Deep Learning}, Ch. 12 (Optimization Algorithms).;%
Deisenroth, Faisal and Ong, \textit{Mathematics for Machine Learning}, Ch. 7.;%
Kingma and Ba, ``Adam: A Method for Stochastic Optimization,'' ICLR 2015.;%
Ruder, ``An overview of gradient descent optimization algorithms,'' arXiv:1609.04747, 2016.%
}

\newcommand{\AMLUnitIVSources}{%
Han, Kamber and Pei, \textit{Data Mining: Concepts and Techniques} (3e), Ch. 3.;%
M\"uller and Guido, \textit{Introduction to Machine Learning with Python}, Ch. 3--4.;%
James, Witten, Hastie, Tibshirani and Taylor, \textit{An Introduction to Statistical Learning with Python}, Ch. 5--6, 12.;%
Scikit-learn documentation (preprocessing, model selection), accessed 2026-08-04.%
}

\newcommand{\AMLUnitVSources}{%
James et al., \textit{An Introduction to Statistical Learning with Python}, Ch. 3--6.;%
Bishop, \textit{Pattern Recognition and Machine Learning}, Ch. 3.;%
Deisenroth, Faisal and Ong, \textit{Mathematics for Machine Learning}, Ch. 9.;%
M\"uller and Guido, \textit{Introduction to Machine Learning with Python}, \S2.3.%
}

\newcommand{\AMLUnitVISources}{%
Han, Kamber and Pei, \textit{Data Mining: Concepts and Techniques} (3e), Ch. 8--9.;%
James et al., \textit{An Introduction to Statistical Learning with Python}, Ch. 8--9.;%
Bishop, \textit{Pattern Recognition and Machine Learning}, Ch. 4--5, 7, 14.;%
Zhang et al., \textit{Dive into Deep Learning}, Ch. 5.%
}

\newcommand{\AMLUnitVIISources}{%
Han, Kamber and Pei, \textit{Data Mining: Concepts and Techniques} (3e), Ch. 10.;%
James et al., \textit{An Introduction to Statistical Learning with Python}, Ch. 12.;%
M\"uller and Guido, \textit{Introduction to Machine Learning with Python}, \S3.5.;%
Ester, Kriegel, Sander and Xu, ``A density-based algorithm for discovering clusters (DBSCAN),'' KDD 1996.%
}

\newcommand{\AMLLabSources}{%
M\"uller and Guido, \textit{Introduction to Machine Learning with Python}, companion notebooks.;%
McKinney, \textit{Python for Data Analysis} (3e), O'Reilly, 2022.;%
Scikit-learn user guide, accessed 2026-08-04.;%
Matplotlib and Seaborn documentation, accessed 2026-08-04.%
}
%
  }{%
    % Question papers live at `private/<family>/<instrument>/`.
    \IfFileExists{../../../public/_shared/aml-sources.tex}{%
      % Source strings for Applied Machine Learning (CSAI2017P) — used by slides + notes.
% Prescribed texts and reference books are taken from the course syllabus.
% Support texts are the author-hosted open-access editions:
%   statlearning.com            (ISL with Python)
%   probml.github.io/pml-book   (Murphy, Probabilistic Machine Learning)
%   d2l.ai                      (Dive into Deep Learning)
%   mml-book.github.io          (Mathematics for Machine Learning)

\newcommand{\AMLRepoURL}{https://github.com/tali7c/Applied-Machine-Learning}

% --- prescribed -------------------------------------------------------------
\newcommand{\AMLPrescribed}{%
M\"uller and Guido, \textit{Introduction to Machine Learning with Python}, O'Reilly, 2016.;%
Bishop, \textit{Pattern Recognition and Machine Learning}, Springer, 2006.;%
Murphy, \textit{Machine Learning: A Probabilistic Perspective}, MIT Press, 2012.;%
Han, Kamber and Pei, \textit{Data Mining: Concepts and Techniques} (3e), Morgan Kaufmann, 2012.%
}

% --- unit-wise ---------------------------------------------------------------
\newcommand{\AMLUnitISources}{%
M\"uller and Guido, \textit{Introduction to Machine Learning with Python}, Ch. 1.;%
Bishop, \textit{Pattern Recognition and Machine Learning}, Ch. 1.;%
Murphy, \textit{Probabilistic Machine Learning: An Introduction}, MIT Press, 2022, Ch. 1.;%
Mitchell, \textit{Machine Learning}, McGraw Hill, 1997, Ch. 1.;%
Scikit-learn documentation, accessed 2026-08-04.%
}

\newcommand{\AMLUnitIISources}{%
Bishop, \textit{Pattern Recognition and Machine Learning}, \S1.5, \S4.3, \S7.1.;%
Zhang, Lipton, Li and Smola, \textit{Dive into Deep Learning}, \S3.4.;%
Shalev-Shwartz and Ben-David, \textit{Understanding Machine Learning}, Ch. 15.;%
Murphy, \textit{Probabilistic Machine Learning: An Introduction}, Ch. 5.%
}

\newcommand{\AMLUnitIIISources}{%
Zhang, Lipton, Li and Smola, \textit{Dive into Deep Learning}, Ch. 12 (Optimization Algorithms).;%
Deisenroth, Faisal and Ong, \textit{Mathematics for Machine Learning}, Ch. 7.;%
Kingma and Ba, ``Adam: A Method for Stochastic Optimization,'' ICLR 2015.;%
Ruder, ``An overview of gradient descent optimization algorithms,'' arXiv:1609.04747, 2016.%
}

\newcommand{\AMLUnitIVSources}{%
Han, Kamber and Pei, \textit{Data Mining: Concepts and Techniques} (3e), Ch. 3.;%
M\"uller and Guido, \textit{Introduction to Machine Learning with Python}, Ch. 3--4.;%
James, Witten, Hastie, Tibshirani and Taylor, \textit{An Introduction to Statistical Learning with Python}, Ch. 5--6, 12.;%
Scikit-learn documentation (preprocessing, model selection), accessed 2026-08-04.%
}

\newcommand{\AMLUnitVSources}{%
James et al., \textit{An Introduction to Statistical Learning with Python}, Ch. 3--6.;%
Bishop, \textit{Pattern Recognition and Machine Learning}, Ch. 3.;%
Deisenroth, Faisal and Ong, \textit{Mathematics for Machine Learning}, Ch. 9.;%
M\"uller and Guido, \textit{Introduction to Machine Learning with Python}, \S2.3.%
}

\newcommand{\AMLUnitVISources}{%
Han, Kamber and Pei, \textit{Data Mining: Concepts and Techniques} (3e), Ch. 8--9.;%
James et al., \textit{An Introduction to Statistical Learning with Python}, Ch. 8--9.;%
Bishop, \textit{Pattern Recognition and Machine Learning}, Ch. 4--5, 7, 14.;%
Zhang et al., \textit{Dive into Deep Learning}, Ch. 5.%
}

\newcommand{\AMLUnitVIISources}{%
Han, Kamber and Pei, \textit{Data Mining: Concepts and Techniques} (3e), Ch. 10.;%
James et al., \textit{An Introduction to Statistical Learning with Python}, Ch. 12.;%
M\"uller and Guido, \textit{Introduction to Machine Learning with Python}, \S3.5.;%
Ester, Kriegel, Sander and Xu, ``A density-based algorithm for discovering clusters (DBSCAN),'' KDD 1996.%
}

\newcommand{\AMLLabSources}{%
M\"uller and Guido, \textit{Introduction to Machine Learning with Python}, companion notebooks.;%
McKinney, \textit{Python for Data Analysis} (3e), O'Reilly, 2022.;%
Scikit-learn user guide, accessed 2026-08-04.;%
Matplotlib and Seaborn documentation, accessed 2026-08-04.%
}
%
    }{%
      % Marking schemes live at `private/solutions/<family>/<id>/latex/`.
      \IfFileExists{../../../../../public/_shared/aml-sources.tex}{%
        % Source strings for Applied Machine Learning (CSAI2017P) — used by slides + notes.
% Prescribed texts and reference books are taken from the course syllabus.
% Support texts are the author-hosted open-access editions:
%   statlearning.com            (ISL with Python)
%   probml.github.io/pml-book   (Murphy, Probabilistic Machine Learning)
%   d2l.ai                      (Dive into Deep Learning)
%   mml-book.github.io          (Mathematics for Machine Learning)

\newcommand{\AMLRepoURL}{https://github.com/tali7c/Applied-Machine-Learning}

% --- prescribed -------------------------------------------------------------
\newcommand{\AMLPrescribed}{%
M\"uller and Guido, \textit{Introduction to Machine Learning with Python}, O'Reilly, 2016.;%
Bishop, \textit{Pattern Recognition and Machine Learning}, Springer, 2006.;%
Murphy, \textit{Machine Learning: A Probabilistic Perspective}, MIT Press, 2012.;%
Han, Kamber and Pei, \textit{Data Mining: Concepts and Techniques} (3e), Morgan Kaufmann, 2012.%
}

% --- unit-wise ---------------------------------------------------------------
\newcommand{\AMLUnitISources}{%
M\"uller and Guido, \textit{Introduction to Machine Learning with Python}, Ch. 1.;%
Bishop, \textit{Pattern Recognition and Machine Learning}, Ch. 1.;%
Murphy, \textit{Probabilistic Machine Learning: An Introduction}, MIT Press, 2022, Ch. 1.;%
Mitchell, \textit{Machine Learning}, McGraw Hill, 1997, Ch. 1.;%
Scikit-learn documentation, accessed 2026-08-04.%
}

\newcommand{\AMLUnitIISources}{%
Bishop, \textit{Pattern Recognition and Machine Learning}, \S1.5, \S4.3, \S7.1.;%
Zhang, Lipton, Li and Smola, \textit{Dive into Deep Learning}, \S3.4.;%
Shalev-Shwartz and Ben-David, \textit{Understanding Machine Learning}, Ch. 15.;%
Murphy, \textit{Probabilistic Machine Learning: An Introduction}, Ch. 5.%
}

\newcommand{\AMLUnitIIISources}{%
Zhang, Lipton, Li and Smola, \textit{Dive into Deep Learning}, Ch. 12 (Optimization Algorithms).;%
Deisenroth, Faisal and Ong, \textit{Mathematics for Machine Learning}, Ch. 7.;%
Kingma and Ba, ``Adam: A Method for Stochastic Optimization,'' ICLR 2015.;%
Ruder, ``An overview of gradient descent optimization algorithms,'' arXiv:1609.04747, 2016.%
}

\newcommand{\AMLUnitIVSources}{%
Han, Kamber and Pei, \textit{Data Mining: Concepts and Techniques} (3e), Ch. 3.;%
M\"uller and Guido, \textit{Introduction to Machine Learning with Python}, Ch. 3--4.;%
James, Witten, Hastie, Tibshirani and Taylor, \textit{An Introduction to Statistical Learning with Python}, Ch. 5--6, 12.;%
Scikit-learn documentation (preprocessing, model selection), accessed 2026-08-04.%
}

\newcommand{\AMLUnitVSources}{%
James et al., \textit{An Introduction to Statistical Learning with Python}, Ch. 3--6.;%
Bishop, \textit{Pattern Recognition and Machine Learning}, Ch. 3.;%
Deisenroth, Faisal and Ong, \textit{Mathematics for Machine Learning}, Ch. 9.;%
M\"uller and Guido, \textit{Introduction to Machine Learning with Python}, \S2.3.%
}

\newcommand{\AMLUnitVISources}{%
Han, Kamber and Pei, \textit{Data Mining: Concepts and Techniques} (3e), Ch. 8--9.;%
James et al., \textit{An Introduction to Statistical Learning with Python}, Ch. 8--9.;%
Bishop, \textit{Pattern Recognition and Machine Learning}, Ch. 4--5, 7, 14.;%
Zhang et al., \textit{Dive into Deep Learning}, Ch. 5.%
}

\newcommand{\AMLUnitVIISources}{%
Han, Kamber and Pei, \textit{Data Mining: Concepts and Techniques} (3e), Ch. 10.;%
James et al., \textit{An Introduction to Statistical Learning with Python}, Ch. 12.;%
M\"uller and Guido, \textit{Introduction to Machine Learning with Python}, \S3.5.;%
Ester, Kriegel, Sander and Xu, ``A density-based algorithm for discovering clusters (DBSCAN),'' KDD 1996.%
}

\newcommand{\AMLLabSources}{%
M\"uller and Guido, \textit{Introduction to Machine Learning with Python}, companion notebooks.;%
McKinney, \textit{Python for Data Analysis} (3e), O'Reilly, 2022.;%
Scikit-learn user guide, accessed 2026-08-04.;%
Matplotlib and Seaborn documentation, accessed 2026-08-04.%
}
%
      }{%
        % Source strings for Applied Machine Learning (CSAI2017P) — used by slides + notes.
% Prescribed texts and reference books are taken from the course syllabus.
% Support texts are the author-hosted open-access editions:
%   statlearning.com            (ISL with Python)
%   probml.github.io/pml-book   (Murphy, Probabilistic Machine Learning)
%   d2l.ai                      (Dive into Deep Learning)
%   mml-book.github.io          (Mathematics for Machine Learning)

\newcommand{\AMLRepoURL}{https://github.com/tali7c/Applied-Machine-Learning}

% --- prescribed -------------------------------------------------------------
\newcommand{\AMLPrescribed}{%
M\"uller and Guido, \textit{Introduction to Machine Learning with Python}, O'Reilly, 2016.;%
Bishop, \textit{Pattern Recognition and Machine Learning}, Springer, 2006.;%
Murphy, \textit{Machine Learning: A Probabilistic Perspective}, MIT Press, 2012.;%
Han, Kamber and Pei, \textit{Data Mining: Concepts and Techniques} (3e), Morgan Kaufmann, 2012.%
}

% --- unit-wise ---------------------------------------------------------------
\newcommand{\AMLUnitISources}{%
M\"uller and Guido, \textit{Introduction to Machine Learning with Python}, Ch. 1.;%
Bishop, \textit{Pattern Recognition and Machine Learning}, Ch. 1.;%
Murphy, \textit{Probabilistic Machine Learning: An Introduction}, MIT Press, 2022, Ch. 1.;%
Mitchell, \textit{Machine Learning}, McGraw Hill, 1997, Ch. 1.;%
Scikit-learn documentation, accessed 2026-08-04.%
}

\newcommand{\AMLUnitIISources}{%
Bishop, \textit{Pattern Recognition and Machine Learning}, \S1.5, \S4.3, \S7.1.;%
Zhang, Lipton, Li and Smola, \textit{Dive into Deep Learning}, \S3.4.;%
Shalev-Shwartz and Ben-David, \textit{Understanding Machine Learning}, Ch. 15.;%
Murphy, \textit{Probabilistic Machine Learning: An Introduction}, Ch. 5.%
}

\newcommand{\AMLUnitIIISources}{%
Zhang, Lipton, Li and Smola, \textit{Dive into Deep Learning}, Ch. 12 (Optimization Algorithms).;%
Deisenroth, Faisal and Ong, \textit{Mathematics for Machine Learning}, Ch. 7.;%
Kingma and Ba, ``Adam: A Method for Stochastic Optimization,'' ICLR 2015.;%
Ruder, ``An overview of gradient descent optimization algorithms,'' arXiv:1609.04747, 2016.%
}

\newcommand{\AMLUnitIVSources}{%
Han, Kamber and Pei, \textit{Data Mining: Concepts and Techniques} (3e), Ch. 3.;%
M\"uller and Guido, \textit{Introduction to Machine Learning with Python}, Ch. 3--4.;%
James, Witten, Hastie, Tibshirani and Taylor, \textit{An Introduction to Statistical Learning with Python}, Ch. 5--6, 12.;%
Scikit-learn documentation (preprocessing, model selection), accessed 2026-08-04.%
}

\newcommand{\AMLUnitVSources}{%
James et al., \textit{An Introduction to Statistical Learning with Python}, Ch. 3--6.;%
Bishop, \textit{Pattern Recognition and Machine Learning}, Ch. 3.;%
Deisenroth, Faisal and Ong, \textit{Mathematics for Machine Learning}, Ch. 9.;%
M\"uller and Guido, \textit{Introduction to Machine Learning with Python}, \S2.3.%
}

\newcommand{\AMLUnitVISources}{%
Han, Kamber and Pei, \textit{Data Mining: Concepts and Techniques} (3e), Ch. 8--9.;%
James et al., \textit{An Introduction to Statistical Learning with Python}, Ch. 8--9.;%
Bishop, \textit{Pattern Recognition and Machine Learning}, Ch. 4--5, 7, 14.;%
Zhang et al., \textit{Dive into Deep Learning}, Ch. 5.%
}

\newcommand{\AMLUnitVIISources}{%
Han, Kamber and Pei, \textit{Data Mining: Concepts and Techniques} (3e), Ch. 10.;%
James et al., \textit{An Introduction to Statistical Learning with Python}, Ch. 12.;%
M\"uller and Guido, \textit{Introduction to Machine Learning with Python}, \S3.5.;%
Ester, Kriegel, Sander and Xu, ``A density-based algorithm for discovering clusters (DBSCAN),'' KDD 1996.%
}

\newcommand{\AMLLabSources}{%
M\"uller and Guido, \textit{Introduction to Machine Learning with Python}, companion notebooks.;%
McKinney, \textit{Python for Data Analysis} (3e), O'Reilly, 2022.;%
Scikit-learn user guide, accessed 2026-08-04.;%
Matplotlib and Seaborn documentation, accessed 2026-08-04.%
}
%
      }%
    }%
  }%
}
%
  }%
}

\usepackage{fancyhdr}
\usepackage{tcolorbox}
\tcbuselibrary{skins,breakable}

\usepackage{array}
\usepackage[htt]{hyphenat}   % allow \texttt{...} to hyphenate, else long
                             % identifiers overflow the measure

\hypersetup{colorlinks=true, linkcolor=blue, urlcolor=blue, breaklinks=true}
\setlist{itemsep=0.35em, topsep=0.35em}
\sloppy
\emergencystretch=2em

\providecommand{\CourseCode}{CSAI2017P: Applied Machine Learning (Lab)}
\providecommand{\AssignmentName}{Lab Assignment}

\pagestyle{fancy}
\fancyhf{}
\lhead{\small\CourseCode}
\rhead{\small\AssignmentName}
\cfoot{\thepage}
\setlength{\headheight}{14pt}

% Per-question mark tag, right aligned.
\newcommand{\QMarks}[1]{\hfill \textbf{[#1 marks]}}

% Title block for a brief.
\newcommand{\LAtitle}[4]{%
  \begin{center}
    {\LARGE \textbf{#1}}\\[0.25em]
    {\large \CourseCode}\\[0.25em]
    \normalsize #2
  \end{center}
  \noindent
  \textbf{Instructor:} Dr.\ Tofik Ali \hfill \textbf{Due:} #3\\
  \textbf{Student Name:} \rule{5cm}{0.4pt} \hfill \textbf{SAP ID:} \rule{3.5cm}{0.4pt}\\
  \textbf{Batch:} \rule{2.5cm}{0.4pt} \hfill \textbf{Lab quiz:} #4
  \vspace{0.6em}\hrule\vspace{0.8em}}

% Title block for a marking scheme (no student fields).
\newcommand{\MStitle}[3]{%
  \begin{center}
    {\LARGE \textbf{#1}}\\[0.25em]
    {\large \CourseCode}\\[0.25em]
    \normalsize #2\\[0.4em]
    {\color{red}\textbf{MARKING SCHEME --- NOT FOR CIRCULATION}}
  \end{center}
  \noindent \textbf{Instructor:} Dr.\ Tofik Ali \hfill \textbf{Assignment due:} #3
  \vspace{0.6em}\hrule\vspace{0.8em}}

\newtcolorbox{infobox}[1][Note]{
  breakable, enhanced, colback=gray!6, colframe=gray!55,
  fonttitle=\bfseries\small, title={#1}, boxrule=0.7pt, arc=2pt,
  left=6pt,right=6pt,top=4pt,bottom=4pt}

\newtcolorbox{answerbox}[1][Expected answer]{
  breakable, enhanced, colback=green!4, colframe=green!35!black,
  fonttitle=\bfseries\small, title={#1}, boxrule=0.7pt, arc=2pt,
  left=6pt,right=6pt,top=4pt,bottom=4pt}


\begin{document}

\LAtitle{Theory Assignment 1: Losses, Optimisers, Preprocessing, Validation and Regression}{Units I--IV and Unit V (part) \textbullet\ Lecture packages L1--L10 \textbullet\ CO1, CO2, CO3}{at AS1, announced in class}{AT1, in class at AS1}

\section*{What this assignment is for}

This is the first of two theory assignments, and it is deliberately built as
your \textbf{mid-semester preparation}. Sections A to C make you rehearse
Units I--IV \emph{and the first three lectures of Unit V} in the three ways an
examination asks for material: explain it, compute it by hand, and apply it in
code. Section D is a full mock paper, to be attempted under examination
conditions.

\begin{warnbox}[The mid-semester syllabus is not confirmed yet]
The examination cell has not issued the notice. \textbf{Units I--IV plus Unit V as far as
L10 is an expectation, not an announcement} --- a well-founded one, because a
paper can only examine what has been taught, and lecture packages L1--L10 are
what the calendar delivers before the mid-semester point.
\textbf{Unit V is covered only to L10:} regression foundations and least
squares (L8), maximum likelihood and model adequacy (L9), and logistic
regression (L10). Multiple linear regression (L11) and regularization (L12) are
\emph{not} examinable here. I will tell you as soon as anything is confirmed.

Nothing about that uncertainty should change what you do now. Everything in
this assignment is examinable material whatever the mid-semester notice says
--- if not in November then in the end-semester paper, which covers all seven
units. The \emph{Mid-Semester Question Bank} released alongside this brief is
organised unit by unit precisely so that it still works if the coverage turns
out narrower or wider than expected; its opening pages explain what to do in
each case.
\end{warnbox}

\begin{infobox}[How the marks work]
The assignment is marked out of \textbf{50 raw}, scaled to the \textbf{1.0}
submission mark. The in-class assessment test \textbf{AT1} that follows at AS1
is worth \textbf{1.5} marks --- more than the submission itself. AT1 asks you
about \emph{your own} submitted work: where a number came from, why you chose
one option over another, what a fragment of your own code does. If you did the
assignment, AT1 takes a few lines per answer. If you did not, no amount of
polish on the submission will save it. This is the design, not an ambush.
\end{infobox}

\begin{center}\begin{tabular}{@{}p{0.8cm}p{5.2cm}p{5.0cm}r@{}}\toprule
\textbf{Sec.} & \textbf{What it tests} & \textbf{Choice rule} & \textbf{Marks}\\
\midrule
A & Concepts, Units I--V(pt) & Attempt any 5 of \textbf{16} (2 marks each) & 10\\
B & Numericals worked by hand & Attempt any 2 of \textbf{12} (5 marks each) & 10\\
C & Pipeline, and a regression fit & Compulsory & 12\\
D & Mock mid-semester paper (out of 35, scaled) & Compulsory & 18\\
\midrule
\multicolumn{3}{@{}l}{\textbf{Raw total}} & \textbf{50}\\
\multicolumn{3}{@{}l}{Scaled to the submission mark} & \textbf{1.0}\\
\multicolumn{3}{@{}l}{Assessment test AT1, in class at AS1} & \textbf{1.5}\\
\bottomrule\end{tabular}\end{center}

\noindent\small Course outcomes assessed: \textbf{CO1, CO2, CO3}. Attempting
more than the choice rule allows is not penalised; your best attempts are
marked.\normalsize

\begin{infobox}[Why there are so many questions you do not have to answer]
Sections A and B offer twenty-eight items and ask for seven. The marks are unchanged
and so is the workload --- but the twenty-eight between them cover every part of
Units I--V(pt) that the mid-semester paper can reach, so the ones you skip are
still worth reading, and the ones you choose are yours. Pick the areas you are
least sure of rather than the ones you can already do; the marks are the same
either way and only one of those choices helps you in the examination hall.

The \textbf{Mid-Semester Question Bank} released alongside this brief carries a
further 314 questions across the same syllabus, 63 of them numerical. It is not
assessed. It is the revision resource.
\end{infobox}

\section*{Instructions}

\begin{itemize}[leftmargin=*]
  \item \textbf{Sections A, B and D are handwritten}, on ruled A4, scanned to a
    single PDF. This is on purpose --- the examination is handwritten and under
    time pressure, and practising on a keyboard does not prepare you for that.
  \item \textbf{Section C is a notebook}, all cells executed, outputs visible.
  \item In Section B, \textbf{show every step}. A correct final number with no
    working earns 1 mark out of 5. A wrong final number with correct method
    earns most of the marks.
  \item Report the metric named in the question. A number with no units or
    context earns nothing.
  \item Set \texttt{random\_state} everywhere in Section C. If your results are
    not reproducible they are not results.
  \item Fit every transformer inside a \texttt{Pipeline}. Nothing is fitted on
    data the model will be evaluated on.
\end{itemize}

\subsection*{What to submit}

\begin{itemize}[leftmargin=*]
  \item \texttt{AML\_A1\_<SAPID>\_written.pdf} --- scanned Sections A, B and D.
  \item \texttt{AML\_A1\_<SAPID>.ipynb} --- Section C, executed.
  \item \texttt{AML\_A1\_<SAPID>\_report.pdf} --- 2--3 pages, or a
    \texttt{README.md}, carrying your Section C results table and observations.
  \item Do not upload datasets. Cite the source and include the loading code.
  \item \textbf{Bring the printed or written submission to AS1.} AT1 is
    open-submission and closed everything else.
\end{itemize}

\begin{warnbox}[Get the data before you start]
Section C uses anchor dataset \textbf{A3 (Statlog German Credit)}, fetched with
\texttt{fetch\_openml('credit-g', version=1, as\_frame=True)}. It downloads
once (\textasciitilde0.2\,MB) and then caches. \textbf{Run the fetch at home,
on your own connection.} The campus network may block it. Full details are in
\texttt{Anchor-Datasets.pdf} alongside the lab assignments.
\end{warnbox}

\section*{Academic integrity}

Discussing ideas with your classmates is encouraged and is how most people
learn this material. Submitting someone else's code or text as your own is not.
If you used a tutorial, a textbook chapter or an AI tool substantially, cite it
in your report --- that is honest practice, not an admission of weakness, and
it costs you no marks. What does cost marks is AT1, which you will not be able
to answer about work you did not do.

\vspace{0.5em}\hrule\vspace{0.8em}

\section*{Section A (Concepts) --- attempt any 5 of 16 \quad (5 $\times$ 2 = 10)}

\noindent\small Answer each in 4--6 lines. Marks are for precision, not length.
The sixteen are ordered by unit: \textbf{A1--A2} Unit I, \textbf{A3--A4} Unit
II, \textbf{A5--A6} Unit III, \textbf{A7--A12} Unit IV,
\textbf{A13--A16} Unit V. \textbf{Choose at least
one from each unit} --- five from one corner of the syllabus is not
revision.\normalsize

\begin{enumerate}[label=\textbf{A\arabic*.}, leftmargin=*]

  \item \textbf{Types of learning} \QMarks{2}\\
    Classify each of the following as supervised, unsupervised or reinforcement
    learning, and state in one line what the training signal is in each case.
    \begin{enumerate}[label=(\alph*), leftmargin=*]
      \item Predicting tomorrow's electricity demand from ten years of hourly meter readings.
      \item Grouping 50{,}000 unlabelled customer profiles into segments for a marketing team.
      \item Training a robot arm to stack blocks, where it is scored only on whether the tower stands.
    \end{enumerate}

  \item \textbf{The three libraries} \QMarks{2}\\
    NumPy, Pandas and scikit-learn each own a distinct job in a Python ML
    workflow. State the job of each in one line. Then explain what
    scikit-learn's \texttt{fit}, \texttt{transform} and \texttt{predict} each do,
    and why \texttt{fit\_transform} exists but \texttt{fit\_predict} is rare.

  \item \textbf{Choosing a regression loss} \QMarks{2}\\
    A house-price model is trained on data where about 2\% of the recorded
    prices are data-entry errors an order of magnitude too large. Compare MSE,
    MAE and Huber loss for this problem. Name the one you would use, and give
    the property of the Huber loss that makes it a compromise between the other
    two.

  \item \textbf{Cross-entropy variants} \QMarks{2}\\
    Binary, categorical and sparse categorical cross-entropy are the same idea
    in three packagings. For each, state (i) the number of classes it applies
    to, (ii) the shape and format the \emph{labels} must be in, and (iii) the
    activation the model's final layer should use.

  \item \textbf{Batch size and momentum} \QMarks{2}\\
    Explain the trade-off between full-batch, mini-batch and stochastic
    gradient descent in terms of gradient noise and cost per update. Then write
    the momentum update rule and state, in one line each, what happens when
    $\beta = 0$ and when $\beta$ is close to $1$.

  \item \textbf{Adagrad, RMSprop and Adam} \QMarks{2}\\
    \begin{enumerate}[label=(\alph*), leftmargin=*]
      \item Write the update rule for each of the three.
      \item State the one failure mode of Adagrad on a long training run, and
        name the single change RMSprop makes to fix it.
      \item Name the two things Adam adds on top of RMSprop.
    \end{enumerate}

  \item \textbf{Fitting the scaler} \QMarks{2}\\
    A classmate standardises the whole dataset, then splits it 80/20 and
    reports test accuracy. Name what has gone wrong, explain the mechanism in
    two lines, and state whether the reported accuracy is likely to be
    optimistic or pessimistic --- with a reason.

  \item \textbf{Imbalance} \QMarks{2}\\
    A fraud dataset is 2\% positive. State one advantage and one danger of
    random oversampling, and the same for random undersampling. Then state
    exactly where in a $k$-fold cross-validation loop the resampling must
    happen, and why it must not happen before the loop.

  \item \textbf{Masking} \QMarks{2}\\
    A column of 12 sensor readings contains two obviously wrong values. The
    $|z| > 3$ rule flags neither of them; the $1.5\times$IQR fence flags both.
    \begin{enumerate}[label=(\alph*), leftmargin=*]
      \item Name the effect and explain the mechanism in two lines.
      \item State the breakdown point of each of the two rules.
      \item Name a third rule with a higher breakdown point than either.
    \end{enumerate}

  \item \textbf{The four scalers} \QMarks{2}\\
    Name \texttt{StandardScaler}, \texttt{MinMaxScaler}, \texttt{RobustScaler}
    and \texttt{MaxAbsScaler}, write each formula, and give the output range
    each guarantees. State which one a single typing error destroys completely,
    and why that one is the most dangerous despite producing output that looks
    perfectly correct.

  \item \textbf{Encoding} \QMarks{2}\\
    A \texttt{colour} column holds \texttt{red}, \texttt{green}, \texttt{blue}.
    A classmate applies \texttt{OrdinalEncoder} because ``it makes fewer
    columns''.
    \begin{enumerate}[label=(\alph*), leftmargin=*]
      \item State the assertion this makes about the variable that is not true.
      \item State which model families believe that assertion and which ignore
        it.
      \item Give the encoding that should have been used and the number of
        columns it produces, with and without \texttt{drop='first'}.
    \end{enumerate}

  \item \textbf{PCA} \QMarks{2}\\
    ``PCA reduces dimensionality by keeping the most important features.''
    State precisely what is wrong with that sentence, say what PCA actually
    keeps, and give one reason a hospital wanting to order fewer blood tests
    would not be helped by it at all.

  \item \textbf{The least-squares line} \QMarks{2}\\
    \begin{enumerate}[label=(\alph*), leftmargin=*]
      \item State exactly which deviations least squares minimises --- be
        precise about the word ``vertical''.
      \item Give the two normal equations, and the estimators $b_1$ and $b_0$
        that follow.
      \item State the four properties of the fitted line, and say which one
        requires the model to have an intercept.
    \end{enumerate}

  \item \textbf{Maximum likelihood} \QMarks{2}\\
    \begin{enumerate}[label=(\alph*), leftmargin=*]
      \item State the noise assumption under which least squares \emph{is} the
        maximum likelihood estimator, in one sentence.
      \item Explain why the value of $\sigma^2$ cannot affect which
        $\beta_0, \beta_1$ maximise the likelihood.
      \item State what maximum likelihood minimises instead if the noise is
        Laplace rather than Gaussian.
    \end{enumerate}

  \item \textbf{$R^2$ and what it hides} \QMarks{2}\\
    A classmate reports $R^2 = 0.87$ on the training set and concludes the
    model is good.
    \begin{enumerate}[label=(\alph*), leftmargin=*]
      \item Give two distinct things wrong with that inference.
      \item State why $R^2$ can never fall when a predictor is added, in one
        sentence.
      \item State the one plot that should accompany every reported $R^2$, and
        what it would reveal that the number cannot.
    \end{enumerate}

  \item \textbf{Odds, log-odds and the sigmoid} \QMarks{2}\\
    \begin{enumerate}[label=(\alph*), leftmargin=*]
      \item Give the range of a probability, of the odds, and of the log-odds.
        State which one a linear expression $w^\top x + b$ can safely produce,
        and why that settles where the model is built.
      \item Invert $\ln\!\big(p/(1-p)\big) = z$ to obtain $p = \sigma(z)$,
        showing the steps.
      \item State what $e^{\beta_j}$ is a change in. \emph{Answering
        ``probability'' scores zero for this part.}
    \end{enumerate}

\end{enumerate}

\section*{Section B (Numericals) --- attempt any 2 of 12 \quad (2 $\times$ 5 = 10)}

\noindent\small Work by hand. Show every step. You may check your answer with
Python afterwards, but the working is what is marked. Give final values to four
decimal places.\normalsize

\begin{enumerate}[label=\textbf{B\arabic*.}, leftmargin=*]

  \item \textbf{Three losses on one set of errors} \QMarks{5}\\
    A regression model produces these predictions:
    \begin{center}\begin{tabular}{@{}lrrrrr@{}}\toprule
    $y$ & 3.0 & $-0.5$ & 2.0 & 7.0 & 4.5\\
    $\hat{y}$ & 2.5 & 0.0 & 2.0 & 8.0 & 1.5\\
    \bottomrule\end{tabular}\end{center}
    \begin{enumerate}[label=(\alph*), leftmargin=*]
      \item Compute the MSE and the MAE. \QMarks{2}
      \item Compute the Huber loss with $\delta = 1$, showing which branch of
        the piecewise definition each of the five points falls into. \QMarks{2}
      \item The fifth point has error 3. Compute what \emph{fraction} of the
        total MSE and of the total Huber loss that single point contributes,
        and say in two lines what this tells you about robustness. \QMarks{1}
    \end{enumerate}

  \item \textbf{Gradient descent, with and without momentum} \QMarks{5}\\
    Fit $\hat{y} = wx$ (no intercept) by minimising
    $L(w) = \frac{1}{n}\sum_i (y_i - wx_i)^2$ on the three points
    $(1,2)$, $(2,4)$, $(3,5)$. Start at $w_0 = 0$ with learning rate
    $\eta = 0.1$.
    \begin{enumerate}[label=(\alph*), leftmargin=*]
      \item Derive $\mathrm{d}L/\mathrm{d}w$ and evaluate it at $w_0$. \QMarks{1}
      \item Perform two full-batch gradient descent updates. Report $w_1$ and $w_2$. \QMarks{1.5}
      \item The least-squares optimum is $w^{*} = \sum x_iy_i / \sum x_i^2$.
        Compute it and state how close $w_2$ is. \QMarks{1}
      \item Now redo the same two steps with momentum, $v_t = \beta v_{t-1} + g_t$
        and $w_t = w_{t-1} - \eta v_t$, with $\beta = 0.9$ and $v_0 = 0$.
        Report $w_1$ and $w_2$, and explain in three lines what momentum has
        done here relative to $w^{*}$. \QMarks{1.5}
    \end{enumerate}

  \item \textbf{Adam against RMSprop against SGD} \QMarks{5}\\
    A single parameter starts at $w_0 = 1.0$. Two successive mini-batches give
    gradients $g_1 = 0.5$ and $g_2 = 0.3$. Use $\eta = 0.1$,
    $\beta_1 = 0.9$, $\beta_2 = 0.999$, $\rho = 0.9$, $\varepsilon = 10^{-8}$.
    \begin{enumerate}[label=(\alph*), leftmargin=*]
      \item Two steps of plain SGD. \QMarks{0.5}
      \item Two steps of Adam. Tabulate $m_t$, $v_t$, $\hat{m}_t$, $\hat{v}_t$
        and $w_t$ at each step. \QMarks{2.5}
      \item Two steps of RMSprop, $v_t = \rho v_{t-1} + (1-\rho)g_t^2$. \QMarks{1}
      \item Adam's first step is almost exactly $-\eta$ regardless of the size
        of $g_1$. Show why from your numbers, and explain in three lines what
        the bias-correction terms $1-\beta_1^t$ and $1-\beta_2^t$ are for.
        \QMarks{1}
    \end{enumerate}

  \item \textbf{Outliers, scaling and the masking effect} \QMarks{5}\\
    A feature column holds the values
    $\;12,\;15,\;14,\;10,\;13,\;60,\;11,\;14$.
    \begin{enumerate}[label=(\alph*), leftmargin=*]
      \item Compute the mean and the population standard deviation
        ($\text{ddof}=0$, as \texttt{StandardScaler} uses). \QMarks{1}
      \item Compute $Q_1$, $Q_3$, the IQR and the $1.5 \times \text{IQR}$
        fences. Which values are flagged as outliers? \QMarks{1.5}
      \item Compute the $z$-score of the value 60. Using the common
        $|z| > 3$ rule, is it flagged? \QMarks{1}
      \item You have two methods disagreeing about the same point. Explain the
        mechanism --- what the value 60 does to the statistic used to judge it
        --- and name which of the two rules you would trust here and why.
        \QMarks{1}
      \item Compute the min--max scaled value of 12 with the outlier present.
        State in two lines what this shows about min--max scaling on
        outlier-contaminated data. \QMarks{0.5}
    \end{enumerate}

  \item \textbf{Cross-entropy, three ways} \QMarks{5}\\
    \begin{enumerate}[label=(\alph*), leftmargin=*]
      \item Four emails have labels $y = (1,0,1,0)$ and predicted
        probabilities of class 1 $p = (0.8,\,0.3,\,0.6,\,0.1)$. Write the
        probability each assigns to its \emph{correct} class, then compute the
        binary cross-entropy. \QMarks{2}
      \item The third email's probability now falls to $0.02$. Recompute the
        BCE and state the increase as a percentage. \QMarks{1}
      \item A three-class model outputs
        $P = \left(\begin{smallmatrix} 0.6 & 0.3 & 0.1\\ 0.2 & 0.7 & 0.1\\
        0.1 & 0.3 & 0.6\end{smallmatrix}\right)$ with true classes $0,1,2$.
        Compute the categorical cross-entropy from one-hot labels and the
        sparse categorical cross-entropy from integer labels. \QMarks{1.5}
      \item Your two answers in (c) are identical. Explain why the sparse form
        exists at all, with a memory argument for $K = 10{,}000$ classes and a
        batch of 256. \QMarks{0.5}
    \end{enumerate}

  \item \textbf{Imputation and the shrinkage it hides} \QMarks{5}\\
    Eight students sat a test; two scripts were lost. The six recorded marks
    are $48,\ 55,\ 60,\ 67,\ 72,\ 78$.
    \begin{enumerate}[label=(\alph*), leftmargin=*]
      \item Compute the mean, median and standard deviation of the six observed
        marks. State clearly which divisor you used. \QMarks{1}
      \item Impute the two missing marks with the mean of the observed. Compute
        the mean and the population standard deviation of all eight. \QMarks{1}
      \item Repeat with median imputation. \QMarks{1}
      \item Compute $\sqrt{m/n}$, where $m$ is the number observed. Compare it
        with the ratio of your imputed standard deviation to the observed one,
        and derive why the two agree. \QMarks{1.5}
      \item State in two lines what this means for any standard error or
        correlation computed after mean imputation. \QMarks{0.5}
    \end{enumerate}

  \item \textbf{PCA by hand} \QMarks{5}\\
    Five points: $(1,1)$, $(2,3)$, $(3,2)$, $(4,5)$, $(5,4)$.
    \begin{enumerate}[label=(\alph*), leftmargin=*]
      \item Compute the mean and write the centred matrix. \QMarks{0.5}
      \item Compute the $2\times2$ covariance matrix using the divisor $n-1$.
        \QMarks{1}
      \item Find both eigenvalues. Verify that they sum to the trace, and give
        the proportion of variance each component explains. \QMarks{1.5}
      \item Give the first eigenvector, project the five centred points onto
        it, and verify that the variance of the five scores equals
        $\lambda_1$. \QMarks{1}
      \item Reconstruct the five points from one component and compute the
        total squared reconstruction error. Verify it equals
        $(n-1)\lambda_2$. \QMarks{1}
    \end{enumerate}

  \item \textbf{The confusion matrix and the metric that lies} \QMarks{5}\\
    A screening model on 10{,}000 patients returns $\text{TN} = 9900$,
    $\text{FP} = 60$, $\text{FN} = 30$, $\text{TP} = 10$.
    \begin{enumerate}[label=(\alph*), leftmargin=*]
      \item Compute accuracy, precision, recall, $F_1$ and balanced accuracy.
        \QMarks{2.5}
      \item Compute the accuracy of a model that predicts ``negative'' for
        every patient. \QMarks{0.5}
      \item State which of the two models scores higher on accuracy, and what
        that tells you about accuracy as a metric here. \QMarks{1}
      \item Name the metric you would put in front of a clinical audit
        committee, and justify it in three lines in terms of the cost of a
        false negative. \QMarks{1}
    \end{enumerate}

  \item \textbf{Least squares by hand} \QMarks{5}\\
    Fit a least-squares line to $x = (0,\,1,\,2,\,3,\,4)$,
    $y = (1,\,2,\,4,\,5,\,8)$.
    \begin{enumerate}[label=(\alph*), leftmargin=*]
      \item Tabulate $x_i-\bar x$, $y_i-\bar y$, $(x_i-\bar x)^2$ and
        $(x_i-\bar x)(y_i-\bar y)$; give $S_{xx}$, $S_{xy}$ and $S_{yy}$.
        \QMarks{1.5}
      \item Compute $b_1$ and $b_0$ and write the fitted line. \QMarks{1}
      \item Give the fitted values and residuals, and verify \emph{both}
        $\sum e_i = 0$ and $\sum x_ie_i = 0$. \QMarks{1}
      \item Compute SSE, SSR and SST, verify the decomposition, and check SSR
        a second independent way. \QMarks{1}
      \item Predict $y$ at $x = 6$. \QMarks{0.5}
    \end{enumerate}

  \item \textbf{Which line is the least-squares line?} \QMarks{5}\\
    Four candidates are proposed for the data of B9.
    \begin{center}\begin{tabular}{@{}llll@{}}\toprule
    A: $0.5 + 1.7x$ & B: $0.6 + 1.7x$ & C: $0.6 + 1.8x$ & D: $1.0 + 1.5x$\\
    \bottomrule\end{tabular}\end{center}
    \begin{enumerate}[label=(\alph*), leftmargin=*]
      \item For each, compute SSE, $\sum e_i$ and $\sum x_ie_i$. \QMarks{3}
      \item Identify the least-squares line using \emph{both} normal
        equations. \QMarks{1}
      \item One candidate satisfies $\sum e_i = 0$ and is still not the
        answer. Name it, and explain in three lines why the first normal
        equation alone cannot pin down a line. \QMarks{1}
    \end{enumerate}

  \item \textbf{$R^2$, and the two estimates of $\sigma^2$} \QMarks{5}\\
    Using your fit from B9 (SSE $= 1.10$, SST $= 30$, $n = 5$):
    \begin{enumerate}[label=(\alph*), leftmargin=*]
      \item Compute $R^2$ from the definition, then verify it two further
        independent ways. \QMarks{1.5}
      \item Compute $r$ and confirm $r^2 = R^2$. \QMarks{0.5}
      \item Compute $\hat\sigma^2_{\text{ML}} = \text{SSE}/n$, the unbiased
        $s^2 = \text{SSE}/(n-2)$, and $s$. \QMarks{1}
      \item Give the bias factor and the percentage by which the maximum
        likelihood estimate understates $\sigma^2$ in expectation. \QMarks{1}
      \item Explain \emph{why} the divisor is $n-2$. \QMarks{1}\\
        \emph{An answer that says ``because we divide by $n$ instead of
        $n-2$'' restates the formula and earns nothing.}
    \end{enumerate}

  \item \textbf{Logistic regression: one gradient step by hand} \QMarks{5}\\
    Take $X = \begin{pmatrix} 1 & -1\\ 1 & 0\\ 1 & 1\\ 1 & 2\end{pmatrix}$,
    $y = (0,\,0,\,1,\,1)$, $w = (0,\,0)$ and $\eta = 0.4$.
    \begin{enumerate}[label=(\alph*), leftmargin=*]
      \item Compute $z$, then $p$, then the log-likelihood. Express the
        per-row negative log-likelihood as a familiar constant and say what it
        is the baseline for. \QMarks{1}
      \item Compute $p - y$ and the gradient $X^\top(p-y)$, showing both
        components. \QMarks{1.5}
      \item Perform one gradient step and give the new $w$. \QMarks{0.5}
      \item Compute the new $z$, the new $p$ and the new log-likelihood, and
        state the improvement. \QMarks{1}
      \item The intercept component of the gradient is exactly zero. Explain
        why, and state what would have to change for it not to be. \QMarks{1}
    \end{enumerate}

\end{enumerate}

\section*{Section C (Pipeline and regression) --- compulsory \quad (12)}

\noindent\small\emph{Dataset: A3 (Statlog German Credit), 1000 rows, 20 mixed
features, binary target \texttt{class} $\in$ \{good, bad\}. Roughly 70/30
imbalanced.}\normalsize

\begin{enumerate}[label=\textbf{C\arabic*.}, leftmargin=*]

  \item \textbf{Build the pipeline} \QMarks{4}
    \begin{enumerate}[label=(\alph*), leftmargin=*]
      \item Load A3. Report the number of rows, the number of numeric and
        categorical features separately, and the class balance.
      \item Build a \texttt{ColumnTransformer} that median-imputes and
        standardises the numeric columns and mode-imputes and one-hot encodes
        the categorical ones. Use
        \texttt{OneHotEncoder(handle\_unknown='ignore')} and say in one line
        why that argument matters.
      \item Wrap it with \texttt{LogisticRegression(max\_iter=2000)} in a
        \texttt{Pipeline}. Split 80/20 with \texttt{stratify=y},
        \texttt{random\_state=42}.
      \item Report the number of columns \emph{after} encoding, the test
        accuracy and the test ROC--AUC.
      \item Run \texttt{StratifiedKFold(5, shuffle=True, random\_state=42)}
        cross-validation on the training portion and report mean $\pm$ standard
        deviation of ROC--AUC.
    \end{enumerate}

  \item \textbf{Make leakage visible} \QMarks{3}\\
    This part uses \emph{pure noise}, so anything above chance is leakage and
    nothing else.
    \begin{enumerate}[label=(\alph*), leftmargin=*]
      \item Generate \texttt{X = rng.normal(size=(200, 5000))} and
        \texttt{y = rng.randint(0, 2, 200)} with
        \texttt{rng = np.random.RandomState(0)}. There is, by construction, no
        relationship between \texttt{X} and \texttt{y}.
      \item \emph{The wrong way.} Apply \texttt{SelectKBest(f\_classif, k=20)}
        to the \emph{whole} of \texttt{X} and \texttt{y}, then cross-validate a
        logistic regression on the selected 20 columns with
        \texttt{StratifiedKFold(5, shuffle=True, random\_state=0)}. Report the
        mean accuracy.
      \item \emph{The right way.} Put the same \texttt{SelectKBest} inside a
        \texttt{Pipeline} with the classifier and cross-validate the pipeline
        on the full \texttt{X}. Report the mean accuracy.
      \item Explain in 5--8 lines why (b) scores what it scores. Your answer
        must name what the selector saw that it should not have, and say which
        of the two numbers an honest report would quote.
    \end{enumerate}

  \item \textbf{Imbalance and the metric you report} \QMarks{3}\\
    Back to A3, treating \texttt{bad} credit as the positive class.
    \begin{enumerate}[label=(\alph*), leftmargin=*]
      \item Report accuracy, precision, recall and ROC--AUC for your C1
        pipeline on the test set.
      \item Refit with \texttt{class\_weight='balanced'} and report the same
        four numbers.
      \item Present both rows in one table. Three of the four numbers move in
        an interesting way and one barely moves. Say which, and explain in
        4--6 lines why ROC--AUC behaves differently from accuracy under a
        change of class weighting.
      \item A bank asks which of the two models to deploy. Answer in three
        lines, and your answer must refer to the \emph{cost} of the two kinds
        of error, not only to the numbers.
    \end{enumerate}

  \item \textbf{Regression, by hand and by library} \QMarks{2}\\
    \emph{No download needed --- use \texttt{sklearn.datasets.load\_diabetes},
    which ships with scikit-learn and works offline.} Take \texttt{bmi} as the
    single predictor and the disease-progression score as the target.
    \begin{enumerate}[label=(\alph*), leftmargin=*]
      \item Split 75/25 with \texttt{random\_state=0}. On the \emph{training}
        rows only, compute $\bar x$, $\bar y$, $S_{xx}$, $S_{xy}$ and hence
        $b_1$ and $b_0$ \textbf{in NumPy, from the formulae} --- no
        \texttt{LinearRegression} in this part. \QMarks{0.75}
      \item Now fit \texttt{LinearRegression()} on the same rows. Print your
        two coefficients beside the library's and report the absolute
        difference. \QMarks{0.5}
      \item Report training and test RMSE, and the test $R^2$. \QMarks{0.5}
      \item Report the RMSE of predicting the training mean for every test row,
        and say in two lines what your $R^2$ means relative to it. \QMarks{0.25}
    \end{enumerate}
    \emph{If (b) does not agree to at least six decimal places, your (a) is
    wrong --- find it before moving on. That agreement is the whole point of
    the part.}

\end{enumerate}

\section*{Section D (Mock mid-semester paper) --- compulsory \quad (18)}

\begin{warnbox}[Attempt this under examination conditions]
\textbf{90 minutes. Closed book. No computer, no calculator beyond arithmetic.}
Sit down, set a timer, and stop when it goes. The paper below is written to the
mid-semester pattern and syllabus \emph{as currently expected} --- see the note
on the first page; neither is confirmed. Marked out of \textbf{35 raw},
scaled to \textbf{18} within this assignment --- the real paper is scaled the same
way to its own weight. Submit your timed script with the rest of the assignment. The answer
key is released after AS1; do not go looking for it first, because the whole
value of this section is in finding out what you cannot yet do while there is
still time to fix it.
\end{warnbox}

\subsection*{Mock paper --- Section A: attempt \emph{all} \quad (5 $\times$ 2 = 10)}

\begin{enumerate}[label=\textbf{\arabic*.}, leftmargin=*]
  \item State exactly which deviations least squares minimises, write the two
    normal equations, and give $b_1$ and $b_0$ in terms of $S_{xx}$ and
    $S_{xy}$. \QMarks{2}
  \item Write the formulae for MSE and MAE. State which is more sensitive to
    outliers and justify it from the formula, not from memory. \QMarks{2}
  \item A 4-class classifier outputs a softmax vector. Your labels are stored
    as the integers $0,1,2,3$. Name the correct loss function, and state what
    you would have to change to use categorical cross-entropy instead. \QMarks{2}
  \item Write the SGD-with-momentum update rule and state in one line what
    $\beta$ controls and what happens as $\beta \to 1$. \QMarks{2}
  \item Define data leakage in one sentence and give two distinct examples of
    preprocessing steps that cause it if performed before the train--test split.
    \QMarks{2}
\end{enumerate}

\subsection*{Mock paper --- Section B: attempt any \emph{3} of 4 \quad (3 $\times$ 5 = 15)}

\begin{enumerate}[label=\textbf{\arabic*.}, leftmargin=*, start=6]
  \item A model produces $y = (10, 12, 15, 9)$ and
    $\hat{y} = (11, 10, 20, 9)$. Compute the MSE, the MAE and the Huber loss
    with $\delta = 2$, showing the branch taken by each point. Comment in two
    lines on which loss the third point dominates. \QMarks{5}
  \item Fit a least-squares line to $x = (0,1,2,3,4)$, $y = (1,2,4,5,8)$.
    Give $S_{xx}$, $S_{xy}$, $b_1$ and $b_0$; the fitted values and residuals;
    and verify \emph{both} $\sum e_i = 0$ and $\sum x_ie_i = 0$. Compute SSE,
    SSR, SST and $R^2$, checking SSR a second independent way. \QMarks{5}
  \item Write the update rules for Adagrad, RMSprop and Adam. State the one
    failure mode of Adagrad on long training runs, name the change RMSprop
    makes to fix it, and name the two things Adam adds on top of RMSprop.
    \QMarks{5}
  \item Take $X$ with rows $(1,-1), (1,0), (1,1), (1,2)$, $y = (0,0,1,1)$,
    $w = (0,0)$ and $\eta = 0.4$.
    \begin{enumerate}[label=(\alph*), leftmargin=*]
      \item Compute $p$ and the log-likelihood, and express the per-row
        negative log-likelihood as a familiar constant. \QMarks{1.5}
      \item Compute the gradient $X^\top(p-y)$, showing both components, and
        perform one gradient step. \QMarks{2}
      \item Give the new $p$ and the new log-likelihood. \QMarks{1}
      \item Explain why the intercept component of the gradient is exactly
        zero here. \QMarks{0.5}
    \end{enumerate}
\end{enumerate}

\subsection*{Mock paper --- Section C: compulsory \quad (1 $\times$ 10 = 10)}

\begin{enumerate}[label=\textbf{\arabic*.}, leftmargin=*, start=10]
  \item A hospital gives you 5000 patient records: 12 numeric features (two of
    them with about 4\% missing values and visible outliers), 3 categorical
    features, and a binary target that is 8\% positive. You must build and
    validate a model, and defend it to a clinical audit committee.
    \begin{enumerate}[label=(\alph*), leftmargin=*]
      \item Set out the complete preprocessing sequence, from raw table to
        model input, in order. For each step state \emph{what} it does and
        \emph{why} it is at that position in the order. \QMarks{3}
      \item State precisely which steps are fitted on training data only, and
        describe what happens inside a single fold of 5-fold stratified
        cross-validation --- fold sizes, number of positives, and which objects
        are refitted. \QMarks{3}
      \item You decide to oversample the minority class. State exactly where in
        the fold this happens, and explain what goes wrong if you oversample
        before the folds are cut. \QMarks{2}
      \item The committee asks for a single headline number. Explain why
        accuracy is a poor choice here, name the metric you would report
        instead, and justify it in terms of the clinical cost of a false
        negative. \QMarks{2}
    \end{enumerate}
\end{enumerate}

\vspace{0.8em}\hrule\vspace{0.6em}
\noindent\small\textbf{End of A1.} The assessment test AT1 follows at AS1 and
is worth 1.5 marks. Bring your submission.\normalsize

\end{document}
